\documentclass[12pt]{article}
\usepackage[ukrainian,shorthands=off]{babel}   % shorthands=off: символ " не активний (потрібно для міток "..." у TikZ всередині \zadtask)
\usepackage{fontspec}
% --- НАЛАШТУВАННЯ ШРИФТІВ ---
% Шрифти Schoolbook-*.otf лежать у корені проєкту. Шлях підбирається автоматично:
% ../ (компіляція з папки теми), ./ (корінь проєкту), ../../ (підпапки теми 28); інакше -- TeX Gyre Schola.
\newcommand{\nmtSchoolbook}[1]{\setmainfont{Schoolbook-Regular.otf}[
    Path = #1,
    BoldFont = Schoolbook-Bold.otf,
    ItalicFont = Schoolbook-Italic.otf,
    BoldItalicFont = Schoolbook-BoldItalic.otf
]}
\IfFileExists{../Schoolbook-Regular.otf}{\nmtSchoolbook{../}}{%
 \IfFileExists{./Schoolbook-Regular.otf}{\nmtSchoolbook{./}}{%
  \IfFileExists{../../Schoolbook-Regular.otf}{\nmtSchoolbook{../../}}{%
   \setmainfont{texgyreschola}[Extension=.otf, UprightFont=*-regular, BoldFont=*-bold, ItalicFont=*-italic, BoldItalicFont=*-bolditalic]}}}
\usepackage[italic]{mathastext}
\usepackage[a4paper,margin=1.8cm,bottom=2cm,top=2cm]{geometry}
\usepackage{amsmath,amssymb}
\usepackage{tikz}
\usepackage{pgfplots}
\pgfplotsset{compat=1.18}
\usepgfplotslibrary{fillbetween}
\usetikzlibrary{calc,patterns,angles,quotes,intersections,babel,3d,shapes.symbols,shapes.geometric,decorations.pathreplacing,shadings,arrows.meta,hobby}
\usepackage{xcolor,array,fancyhdr,enumitem,multicol}
\usepackage{colortbl,multirow,diagbox,needspace}
\usepackage{tcolorbox}
\tcbuselibrary{skins,breakable}
\usepackage[hidelinks]{hyperref}
\usepackage[firstpage=false, text=@pvtr2525, scale=0.6, color=gray!12]{draftwatermark}

% --- АВТОСТИСКАННЯ: рисунок/сітка, ширша за колонку, масштабується до ширини колонки ---
\newsavebox{\nmtfitbox}
\newenvironment{nmtfit}{\begin{lrbox}{\nmtfitbox}}{\end{lrbox}%
  \ifdim\wd\nmtfitbox>\linewidth \resizebox{\linewidth}{!}{\usebox{\nmtfitbox}}\else\usebox{\nmtfitbox}\fi}
\let\nmtIncludegraphicsOrig\includegraphics
\renewcommand{\includegraphics}[2][]{\begin{nmtfit}\nmtIncludegraphicsOrig[#1]{#2}\end{nmtfit}}


\definecolor{mainGreen}{RGB}{34, 120, 64}
\definecolor{yearOrange}{RGB}{220, 100, 30}
\definecolor{yearcolor}{RGB}{220, 100, 30}
\definecolor{titleBg}{RGB}{225, 240, 220}
\definecolor{instrBg}{RGB}{255, 240, 220}
\definecolor{instrBorder}{RGB}{220, 150, 80}
\definecolor{headerblue}{RGB}{0, 102, 204}   % використовується в рисунках старої бази

\setlength{\parindent}{0pt}
\emergencystretch=1.5em

\pagestyle{fancy}
\fancyhf{}
\renewcommand{\headrulewidth}{0.4pt}
\renewcommand{\footrulewidth}{0.4pt}
\fancyhead[C]{\small\color{gray!80} База завдань НМТ 2022--2026 \textendash{} Тема 20}
\fancyhead[R]{\small\color{mainGreen}\textbf{\href{https://t.me/pvtr2525}{@pvtr2525}}}
\fancyfoot[L]{\small\color{gray!80}\href{https://t.me/pvtr2525}{ТГ: @pvtr2525}}
\fancyfoot[C]{\small\color{gray!80}\thepage}
\fancyfoot[R]{\small\color{gray!80}НМТ 2022--2026}

% --- ТАБЛИЦІ ВІДПОВІДЕЙ ---
% ширина клітинки = (ширина рядка - відступи) / 5: на всю сторінку це ~3 см, у minipage поруч із рисунком таблиця стискається
\newlength{\nmtcellw}
\newcommand{\nmtSetCell}{\setlength{\nmtcellw}{\dimexpr(\linewidth-12\tabcolsep-6\arrayrulewidth)/5\relax}}
\newcommand{\answerTable}[5]{%
\par\nopagebreak\vspace{0.15cm}
\noindent\nmtSetCell
\begin{tabular}{|*{5}{>{\centering\arraybackslash}m{\nmtcellw}|}}
\hline
\rule[-0.15cm]{0pt}{0.6cm}\textbf{А} & \rule[-0.15cm]{0pt}{0.6cm}\textbf{Б} & \rule[-0.15cm]{0pt}{0.6cm}\textbf{В} & \rule[-0.15cm]{0pt}{0.6cm}\textbf{Г} & \rule[-0.15cm]{0pt}{0.6cm}\textbf{Д} \\
\hline
\rule[-0.2cm]{0pt}{0.75cm}#1 & \rule[-0.2cm]{0pt}{0.75cm}#2 & \rule[-0.2cm]{0pt}{0.75cm}#3 & \rule[-0.2cm]{0pt}{0.75cm}#4 & \rule[-0.2cm]{0pt}{0.75cm}#5 \\
\hline
\end{tabular}
\par\vspace{0.25cm}
}
\newcommand{\answerTableTall}[5]{%
\par\nopagebreak\vspace{0.15cm}
\noindent\nmtSetCell
\begin{tabular}{|*{5}{>{\centering\arraybackslash}m{\nmtcellw}|}}
\hline
\rule[-0.2cm]{0pt}{0.7cm}\textbf{А} & \rule[-0.2cm]{0pt}{0.7cm}\textbf{Б} & \rule[-0.2cm]{0pt}{0.7cm}\textbf{В} & \rule[-0.2cm]{0pt}{0.7cm}\textbf{Г} & \rule[-0.2cm]{0pt}{0.7cm}\textbf{Д} \\
\hline
\rule[-0.45cm]{0pt}{1.2cm}#1 & \rule[-0.45cm]{0pt}{1.2cm}#2 & \rule[-0.45cm]{0pt}{1.2cm}#3 & \rule[-0.45cm]{0pt}{1.2cm}#4 & \rule[-0.45cm]{0pt}{1.2cm}#5 \\
\hline
\end{tabular}
\par\vspace{0.25cm}
}
\newcommand{\instructionBox}[1]{%
\par\vspace{0.3cm}
\begin{tcolorbox}[colback=instrBg, colframe=instrBorder, boxrule=0.6pt, arc=2pt,
    left=8pt, right=8pt, top=4pt, bottom=4pt]
\centering\small\bfseries #1
\end{tcolorbox}
\par\vspace{0.15cm}
}
\newcommand{\sectionTitle}[1]{%
\par\vspace{0.4cm}
\begin{tcolorbox}[colback=titleBg, colframe=titleBg, boxrule=0pt, arc=8pt,
    halign=center, fontupper=\Large\bfseries\color{mainGreen}]
\hyphenpenalty=10000 \exhyphenpenalty=10000 #1
\end{tcolorbox}
\par\vspace{0.3cm}
}
\newcommand{\task}[2]{\noindent\textbf{#1.}\ \ #2\par\vspace{0.2cm}}
% --- НМТ-СТИЛЬ ПОЛЕ ВІДПОВІДІ ---
\newcommand{\nmtAnswerBox}{%
\par\nopagebreak\vspace{0.2cm}\noindent
Відповідь:\ \
\begin{tabular}{@{}*{4}{|>{\centering\arraybackslash}p{0.8cm}}|@{\hspace{0.1cm}\textbf{,}\hspace{0.1cm}}*{3}{|>{\centering\arraybackslash}p{0.8cm}}|@{}}
\hline
\rule[-0.2cm]{0pt}{0.7cm} & & & & & & \\
\hline
\end{tabular}
\par\vspace{0.3cm}
}
% --- СІТКА ДЛЯ МАТЧИНГУ ---
\newsavebox{\nmtgridbox}
\newcommand{\matchingGrid}{%
    \sbox{\nmtgridbox}{\matchingGridRaw}%
    \ifdim\wd\nmtgridbox>\linewidth \resizebox{\linewidth}{!}{\usebox{\nmtgridbox}}\else\usebox{\nmtgridbox}\fi}
\newcommand{\matchingGridRaw}{%
    \begingroup
    \renewcommand{\arraystretch}{1.15}
    \setlength{\tabcolsep}{4pt}
    \footnotesize
    \begin{tabular}{r|c|c|c|c|c|}
         \multicolumn{1}{c}{} & \multicolumn{1}{c}{\textbf{А}} & \multicolumn{1}{c}{\textbf{Б}} & \multicolumn{1}{c}{\textbf{В}} & \multicolumn{1}{c}{\textbf{Г}} & \multicolumn{1}{c}{\textbf{Д}} \\ \cline{2-6}
         \textbf{1} & \phantom{X} & \phantom{X} & \phantom{X} & \phantom{X} & \phantom{X} \\ \cline{2-6}
         \textbf{2} & & & & & \\ \cline{2-6}
         \textbf{3} & & & & & \\ \cline{2-6}
    \end{tabular}
    \endgroup
}
% --- ДОДАТКОВІ МАКРОСИ ЗБІРНИКА ---
\newcommand{\taskBlock}[1]{%
\par\vspace{0.45cm}
\noindent{\color{mainGreen}\rule{\linewidth}{0.8pt}}\par\vspace{0.12cm}
\noindent{\small\bfseries\color{mainGreen}#1}\par\vspace{0.25cm}
}
\newcounter{zad}
\newcommand{\zadtask}[1]{\stepcounter{zad}\noindent\textbf{\thezad.}\ \ #1\par\nopagebreak\vspace{0.2cm}}
\newcommand{\zadnum}{\stepcounter{zad}\textbf{\thezad.}\ \ }
\newcounter{ans}
\newcommand{\ansitem}[1]{\stepcounter{ans}\noindent\textbf{\theans.}\ #1\par\vspace{0.07cm}}
% мітка року: нерозривна, притиснута праворуч; якщо не вміщається -- праворуч на новому рядку
\newcommand{\nmtyear}[1]{\unskip\nobreak\finalhyphendemerits=0 \hfill\penalty100\hskip1em\hbox{}\nobreak\hfill{\small\color{yearOrange}\mbox{\rule{0pt}{2.3ex}(НМТ~#1)}}}
% --- МАКРОС КОРОТКОГО РОЗВ'ЯЗКУ ---
\newcommand{\solution}[1]{%
\par\vspace{0.12cm}
\begin{tcolorbox}[colback=titleBg!45, colframe=mainGreen!55, boxrule=0.4pt, arc=2pt,
    left=6pt, right=6pt, top=3pt, bottom=3pt, breakable]
\footnotesize\textbf{Короткий розв'язок.}\ #1
\end{tcolorbox}
\par\vspace{0.12cm}
}
% Розв'язки й відповіді (є до завдань НМТ-2026): \showsolutionstrue -- показати, \showsolutionsfalse -- сховати
\newif\ifshowsolutions
\showsolutionsfalse
% --- ЗАГОЛОВКИ З ПУНКТАМИ КЛІКАБЕЛЬНОГО ЗМІСТУ ---
\setcounter{tocdepth}{2}
\newcommand{\chapterTitle}[1]{%
\clearpage\phantomsection\addcontentsline{toc}{section}{#1}%
\sectionTitle{#1}}
\newcommand{\typeTitle}[1]{%
\needspace{6\baselineskip}%
\phantomsection\addcontentsline{toc}{subsection}{\quad #1}%
\par\vspace{0.3cm}
\noindent{\large\bfseries\color{yearOrange}#1}\par\vspace{0.08cm}
\noindent{\color{yearOrange}\rule{\linewidth}{0.4pt}}\par\nopagebreak\vspace{0.2cm}}
\raggedbottom
% усередині одного завдання сторінка не розривається: нерозривний вертикальний проміжок
% і нерозривна межа абзаців (умова ніколи не відділяється від варіантів, рисунка чи сітки)
\newcommand{\nbvspace}[1]{\nopagebreak\vspace{#1}}
\newcommand{\nmtnobreak}{\ifvmode\penalty10000 \fi}
% --- ЄДИНИЙ МАКЕТ ЗАВДАНЬ НА ВІДПОВІДНІСТЬ ---
% мітка в колонці сталої ширини, текст -- у parbox: продовження довгого рядка
% вирівнюється під текстом, а не під міткою; відступ між пунктами однаковий скрізь
\newlength{\nmtMatchLab}\setlength{\nmtMatchLab}{1.6em}
\newlength{\nmtMatchSep}\setlength{\nmtMatchSep}{0.28cm}
\newcommand{\matchHead}[1]{\par\noindent\textit{#1}\par\nopagebreak\vspace{0.2cm}}
\newcommand{\matchItem}[2]{\par\noindent\makebox[\nmtMatchLab][l]{\textbf{#1}}%
\parbox[t]{\dimexpr\linewidth-\nmtMatchLab\relax}{\raggedright #2}\par\nopagebreak\vspace{\nmtMatchSep}}
\newcommand{\ansTheme}[1]{\par\vspace{0.25cm}\noindent{\bfseries\color{mainGreen}#1}\par\vspace{0.12cm}}
\newcommand{\ansType}[1]{\par\vspace{0.1cm}\noindent{\itshape\color{yearOrange}#1}\par\vspace{0.08cm}}

\begin{document}
\begingroup\setcounter{zad}{0}
% --- макроси, перенесені зі старого файлу теми ---
\newcommand{\matchingLayout}[3]{
    \noindent
    \begin{minipage}[t]{0.36\textwidth}\vspace{0pt}\raggedright
       
        #1
    \end{minipage}%
    \hfill
    \begin{minipage}[t]{0.43\textwidth}\vspace{0pt}\raggedright
        
        #2
    \end{minipage}%
    \hfill
    \begin{minipage}[t]{0.19\textwidth}\vspace{0pt}\raggedright
        \vspace{0pt} % Хаки для вирівнювання minipage по верху
        \begin{flushright}
        #3
        \end{flushright}
    \end{minipage}
}

% --- макроси з файлів сесій НМТ-2026 (для рисунків) ---
\renewcommand{\headrulewidth}{0.4pt}
\renewcommand{\footrulewidth}{0.4pt}
\definecolor{gold}{RGB}{240, 190, 40}
\definecolor{silver}{RGB}{170, 170, 175}
\definecolor{bronze}{RGB}{190, 120, 70}
\definecolor{bar1}{RGB}{200, 120, 20}
\definecolor{bar2}{RGB}{255, 220, 0}
\definecolor{bar3}{RGB}{220, 30, 30}
\definecolor{bar4}{RGB}{128, 0, 128}
\definecolor{bar5}{RGB}{0, 160, 220}
\definecolor{bar6}{RGB}{120, 180, 50}
\definecolor{hexcol}{RGB}{200,110,60}
\newcommand{\hexthree}[2]{%
    \draw[hexcol, line width=1.1pt]
        ($(#1,#2)+(0:30)$) -- ($(#1,#2)+(60:30)$) -- ($(#1,#2)+(120:30)$) --
        ($(#1,#2)+(180:30)$) -- ($(#1,#2)+(240:30)$) -- ($(#1,#2)+(300:30)$) -- cycle;
}

% --- сумісність зі старою базою (діє, якщо тема не має власного визначення) ---
\providecommand{\trigStyle}[1]{#1}
\providecommand{\answerTableSmall}[5]{%
\begingroup\setlength{\tabcolsep}{2pt}\small\nmtSetCell
\begin{tabular}{|*{5}{>{\centering\arraybackslash}m{\nmtcellw}|}}
\hline
\rule[-0.2cm]{0pt}{0.6cm}\textbf{А} & \textbf{Б} & \textbf{В} & \textbf{Г} & \textbf{Д} \\
\hline
\rule[-0.4cm]{0pt}{0.9cm}#1 & \rule[-0.4cm]{0pt}{0.9cm}#2 & \rule[-0.4cm]{0pt}{0.9cm}#3 & \rule[-0.4cm]{0pt}{0.9cm}#4 & \rule[-0.4cm]{0pt}{0.9cm}#5 \\
\hline
\end{tabular}\endgroup}
\providecommand{\matchingLayout}[3]{%
    \noindent
    \begin{minipage}[t]{0.36\textwidth}\vspace{0pt}\raggedright
        #1
    \end{minipage}%
    \hfill
    \begin{minipage}[t]{0.43\textwidth}\vspace{0pt}\raggedright
        #2
    \end{minipage}%
    \hfill
    \begin{minipage}[t]{0.19\textwidth}
        \vspace{0pt}
        \begin{flushright}
        #3
        \end{flushright}
    \end{minipage}}
\providecommand{\matchingLayoutBottom}[2]{%
    \noindent
    \begin{minipage}[t]{0.48\textwidth}\vspace{0pt}\raggedright
        #1
    \end{minipage}%
    \hfill
    \begin{minipage}[t]{0.48\textwidth}\vspace{0pt}\raggedright
        #2
    \end{minipage}
    \vspace{0.5cm}
    \noindent
    \matchingGrid}
\providecommand{\answerListVertical}[5]{%
    \par\nopagebreak\vspace{0.2cm}
    \begin{itemize}[itemsep=0.4cm, leftmargin=1.5cm, labelsep=0.5cm]
        \item[\textbf{А}] #1
        \item[\textbf{Б}] #2
        \item[\textbf{В}] #3
        \item[\textbf{Г}] #4
        \item[\textbf{Д}] #5
    \end{itemize}
    \vspace{0.2cm}}

\chapterTitle{Тема 20. Функція та її властивості}
\noindent{\small\color{gray!80} Розділ 3. Функції \quad\textbullet\quad Усього завдань: 63 (НМТ \mbox{2023~--- 17}, \mbox{2024~--- 11}, \mbox{2025~--- 15}, \mbox{2026~--- 20}); \mbox{тестових~--- 29}; \mbox{на відповідність~--- 31}; \mbox{з короткою відповіддю~--- 3}.}\par\vspace{0.2cm}

\typeTitle{Завдання з вибором однієї відповіді (А--Д)}

\begin{samepage}
% === ЗАВДАННЯ 4 (Чверті 1) ===
\noindent\zadnum \begin{minipage}[t]{0.55\textwidth}
Функція $y=f(x)$ визначена на проміжку $(-\infty; +\infty)$ і набуває лише додатних значень. Укажіть \textit{усі} координатні чверті (див. рисунок), у яких розташований графік цієї функції. \nmtyear{2023}
\end{minipage}
\hfill
\begin{minipage}[t]{0.4\textwidth}
    \nopagebreak\nbvspace{-0.3cm}
    \begin{flushright}
    \begin{nmtfit}\begin{tikzpicture}[scale=0.8]
        \draw[->, >=stealth, thick] (-2.5,0) -- (2.5,0) node[below] {$x$};
        \draw[->, >=stealth, thick] (0,-2.5) -- (0,2.5) node[left] {$y$};
        \node[below left] at (0,0) {$O$};

        \node at (1.5,1.5) {I чверть};
        \node at (-1.5,1.5) {II чверть};
        \node at (-1.5,-1.5) {III чверть};
        \node at (1.5,-1.5) {IV чверть};
    \end{tikzpicture}\end{nmtfit}
    \end{flushright}
\end{minipage}

\nmtnobreak
\nopagebreak\nbvspace{0.3cm}
\begin{tabular}{|*{5}{>{\centering\arraybackslash}m{2.8cm}|}}
\hline
\textbf{А} & \textbf{Б} & \textbf{В} & \textbf{Г} & \textbf{Д} \\
\hline
лише в I та IV & лише в II & лише в III та IV & лише в I та II & лише в I \\
\hline
\end{tabular}
\par\penalty-20\vspace{0.25cm}
\end{samepage}

\begin{samepage}
% === ЗАВДАННЯ 5 (Чверті 2) ===
\noindent\zadnum \begin{minipage}[t]{0.55\textwidth}
Функція $y=f(x)$ визначена на проміжку $(0; +\infty)$ і набуває лише додатних значень. Укажіть \textit{усі} координатні чверті (див. рисунок), у яких розташований графік цієї функції. \nmtyear{2023}
\end{minipage}
\hfill
\begin{minipage}[t]{0.4\textwidth}
    \nopagebreak\nbvspace{-0.3cm}
    \begin{flushright}
    \begin{nmtfit}\begin{tikzpicture}[scale=0.8]
        \draw[->, >=stealth, thick] (-2.5,0) -- (2.5,0) node[below] {$x$};
        \draw[->, >=stealth, thick] (0,-2.5) -- (0,2.5) node[left] {$y$};
        \node[below left] at (0,0) {$O$};

        \node at (1.5,1.5) {I чверть};
        \node at (-1.5,1.5) {II чверть};
        \node at (-1.5,-1.5) {III чверть};
        \node at (1.5,-1.5) {IV чверть};
    \end{tikzpicture}\end{nmtfit}
    \end{flushright}
\end{minipage}

\nmtnobreak
\nopagebreak\nbvspace{0.3cm}
\begin{tabular}{|*{5}{>{\centering\arraybackslash}m{2.8cm}|}}
\hline
\textbf{А} & \textbf{Б} & \textbf{В} & \textbf{Г} & \textbf{Д} \\
\hline
лише в I & лише в I та IV & лише в III та IV & лише в I та II & лише в II \\
\hline
\end{tabular}
\par\penalty-20\vspace{0.25cm}
\end{samepage}

\begin{samepage}
% === ЗАВДАННЯ 6 (Чверті 3) ===
\noindent\zadnum \begin{minipage}[t]{0.55\textwidth}
Функція $y=f(x)$ визначена на проміжку $(-\infty; +\infty)$ і набуває лише від’ємних значень. Укажіть \textit{усі} координатні чверті (див. рисунок), у яких розташований графік цієї функції. \nmtyear{2023}
\end{minipage}
\hfill
\begin{minipage}[t]{0.4\textwidth}
    \nopagebreak\nbvspace{-0.3cm}
    \begin{flushright}
    \begin{nmtfit}\begin{tikzpicture}[scale=0.8]
        \draw[->, >=stealth, thick] (-2.5,0) -- (2.5,0) node[below] {$x$};
        \draw[->, >=stealth, thick] (0,-2.5) -- (0,2.5) node[left] {$y$};
        \node[below left] at (0,0) {$O$};

        \node at (1.5,1.5) {I чверть};
        \node at (-1.5,1.5) {II чверть};
        \node at (-1.5,-1.5) {III чверть};
        \node at (1.5,-1.5) {IV чверть};
    \end{tikzpicture}\end{nmtfit}
    \end{flushright}
\end{minipage}

\nmtnobreak
\nopagebreak\nbvspace{0.3cm}
\begin{tabular}{|*{5}{>{\centering\arraybackslash}m{2.8cm}|}}
\hline
\textbf{А} & \textbf{Б} & \textbf{В} & \textbf{Г} & \textbf{Д} \\
\hline
лише в III та IV & лише в IV & лише в III & лише в I та IV & лише в I \\
\hline
\end{tabular}
\par\penalty-20\vspace{0.25cm}
\end{samepage}

\begin{samepage}
\noindent\zadnum \begin{minipage}[t]{0.55\textwidth}
На рисунку зображено графік функції $y=f(x)$, визначеної на проміжку $[-3; 3]$. На якому з наведених проміжків ця функція зростає? \nmtyear{2023}
\end{minipage}
\hfill
\begin{minipage}[t]{0.4\textwidth}
    \nopagebreak\nbvspace{-0.3cm}
    \begin{flushright}
    \begin{nmtfit}\begin{tikzpicture}[scale=0.5]
        \draw[step=1cm,gray!50,very thin] (-3.5,-2.5) grid (3.5,3.5);
        \draw[->, >=stealth, thick] (-3.5,0) -- (3.5,0) node[below] {$x$};
        \draw[->, >=stealth, thick] (0,-2.5) -- (0,3.5) node[left] {$y$};

        \node[below left] at (0,0) {$0$};
        \node[below] at (1,0) {$1$};
        \node[left] at (0,1) {$1$};
        \node[below] at (3,0) {$3$};
        \node[below] at (-3,0) {$-3$};
        \draw (1,0.1) -- (1,-0.1);
        \draw (3,0.1) -- (3,-0.1);
        \draw (-3,0.1) -- (-3,-0.1);
        \draw (0.1,1) -- (-0.1,1);

        % Виправлено: Максимум чітко в точці (1, 3)
        % Проходить через (-3, -2), перетинає вісь десь біля -2, пік (1,3), кінець (3, 1.5)
        \draw[thick] plot [smooth, tension=0.25] coordinates {(-3,-2) (-1.8,0) (1,3) (3,1.5)};

        \fill (-3,-2) circle (3pt);
        \fill (3,1.5) circle (3pt);
        \node[above right, fill=white, inner sep=1.5pt] at (2,2.5) {$y=f(x)$};
    \end{tikzpicture}\end{nmtfit}
    \end{flushright}
\end{minipage}

\nmtnobreak
\nopagebreak\nbvspace{0.3cm}
\answerTable{$[-3; 3]$}{$[-2; 4]$}{$[-2; 3]$}{$[1; 3]$}{$[-3; 1]$}
\par\penalty-20
\end{samepage}

\begin{samepage}
\noindent\zadnum \begin{minipage}[t]{0.55\textwidth}
На рисунку зображено графік функції $y=f(x)$, визначеної на проміжку $[-3; 3]$. Укажіть нуль цієї функції. \nmtyear{2023}
\end{minipage}
\hfill
\begin{minipage}[t]{0.4\textwidth}
    \nopagebreak\nbvspace{-0.3cm}
    \begin{flushright}
    \begin{nmtfit}\begin{tikzpicture}[scale=0.5]
        \draw[step=1cm,gray!50,very thin] (-3.5,-2.5) grid (3.5,3.5);
        \draw[->, >=stealth, thick] (-3.5,0) -- (3.5,0) node[below] {$x$};
        \draw[->, >=stealth, thick] (0,-2.5) -- (0,3.5) node[left] {$y$};

        \node[below left] at (0,0) {$0$};
        \node[below] at (1,0) {$1$};
        \node[left] at (0,1) {$1$};
        \node[below] at (3,0) {$3$};
        \node[below] at (-3,0) {$-3$};
        \draw (1,0.1) -- (1,-0.1);
        \draw (3,0.1) -- (3,-0.1);
        \draw (-3,0.1) -- (-3,-0.1);
        \draw (0.1,1) -- (-0.1,1);

        % Той самий графік, що і в завданні 7
        \draw[thick] plot [smooth, tension=0.6] coordinates {(-3,-2) (-2,0) (1,3) (3,1.5)};

        \fill (-3,-2) circle (3pt);
        \fill (3,1.5) circle (3pt);
        \node[above right, fill=white, inner sep=1.5pt] at (2,2.5) {$y=f(x)$};
    \end{tikzpicture}\end{nmtfit}
    \end{flushright}
\end{minipage}

\nmtnobreak
\nopagebreak\nbvspace{0.3cm}
\answerTable{$0$}{$3$}{$4$}{$-3$}{$-2$}
\par\penalty-20
\end{samepage}

\begin{samepage}
\noindent\zadnum \begin{minipage}[t]{0.55\textwidth}
На рисунку зображено графік функції $y=f(x)$, визначеної на проміжку $[-2; 4]$. Укажіть значення $x_0$, за якого $f(x_0) > 0$. \nmtyear{2023}
\end{minipage}
\hfill
\begin{minipage}[t]{0.4\textwidth}
    \nopagebreak\nbvspace{-0.3cm}
    \begin{flushright}
    \begin{nmtfit}\begin{tikzpicture}[scale=0.5]
        \draw[step=1cm,gray!50,very thin] (-2.5,-2.5) grid (4.5,3.5);
        \draw[->, >=stealth, thick] (-2.5,0) -- (4.5,0) node[below] {$x$};
        \draw[->, >=stealth, thick] (0,-2.5) -- (0,3.5) node[left] {$y$};

        \node[below left] at (0,0) {$0$};
        \node[below] at (1,0) {$1$};
        \node[left] at (0,1) {$1$};
        \node[below] at (4,0) {$4$};
        \node[below] at (-2,0) {$-2$};

        \draw (1,0.1) -- (1,-0.1);
        \draw (4,0.1) -- (4,-0.1);
        \draw (-2,0.1) -- (-2,-0.1);
        \draw (0.1,1) -- (-0.1,1);

        % Виправлено: 
        % Точка (-1, 3)
        % Точка перетину з Oy: (0, 2)
        % Точка перетину з Ox: (1, 0)
        \draw[thick] plot [smooth, tension=0.4] coordinates {(-2, 1.2) (-1, 3) (0, 2) (1, 0) (4, -1.8)};

        \fill (-2,1.2) circle (3pt);
        \fill (4,-1.8) circle (3pt);
        \node[above right, fill=white, inner sep=1.5pt] at (2,1) {$y=f(x)$};
    \end{tikzpicture}\end{nmtfit}
    \end{flushright}
\end{minipage}

\nmtnobreak
\nopagebreak\nbvspace{0.3cm}
\answerTable{$x_0 = 2$}{$x_0 = -1$}{$x_0 = 3$}{$x_0 = 1$}{$x_0 = 4$}
\par\penalty-20
\end{samepage}

\begin{samepage}
\noindent\zadnum \begin{minipage}[t]{0.55\textwidth}
Функція $y=f(x)$ визначена на проміжку $(-\infty; 0)$ і набуває лише від’ємних значень. Укажіть \textit{усі} координатні чверті (див. рисунок), у яких розташований графік цієї функції. \nmtyear{2023}
\end{minipage}
\hfill
\begin{minipage}[t]{0.4\textwidth}
    \nopagebreak\nbvspace{-0.3cm}
    \begin{flushright}
    \begin{nmtfit}\begin{tikzpicture}[scale=0.7]
        \draw[->, >=stealth, thick] (-2.5,0) -- (2.5,0) node[below] {$x$};
        \draw[->, >=stealth, thick] (0,-2.5) -- (0,2.5) node[left] {$y$};
        \node[below left] at (0,0) {$O$};

        \node at (1.9,1.5) {I чверть};
        \node at (-1.9,1.5) {II чверть};
        \node at (-1.9,-1.5) {III чверть};
        \node at (1.9,-1.5) {IV чверть};
    \end{tikzpicture}\end{nmtfit}
    \end{flushright}
\end{minipage}

\nmtnobreak
\nopagebreak\nbvspace{0.3cm}
\answerTable{лише в III}{лише в II та III}{лише в II}{лише в IV}{лише в III та IV}
\par\penalty-20
\end{samepage}

\begin{samepage}
\noindent\zadnum \begin{minipage}[t]{0.55\textwidth}
На рисунку зображено графік функції $y=f(x)$, визначеної на проміжку $[-2; 4]$. Цей графік перетинає вісь $x$ в одній із зазначених точок. Укажіть цю точку. \nmtyear{2023}
\end{minipage}
\hfill
\begin{minipage}[t]{0.4\textwidth}
    \nopagebreak\nbvspace{-0.3cm}
    \begin{flushright}
    \begin{nmtfit}\begin{tikzpicture}[scale=0.5]
        \draw[step=1cm,gray!50,very thin] (-2.5,-1.5) grid (4.5,4.5);
        \draw[->, >=stealth, thick] (-2.5,0) -- (4.5,0) node[below] {$x$};
        \draw[->, >=stealth, thick] (0,-1.5) -- (0,4.5) node[left] {$y$};

        \node[below left] at (0,0) {$0$};
        \node[below] at (1,0) {$1$};
        \node[left] at (0,1) {$1$};
        \node[below] at (4,0) {$4$};
        \node[below] at (-2,0) {$-2$};

        \draw (1,0.1) -- (1,-0.1);
        \draw (4,0.1) -- (4,-0.1);
        \draw (-2,0.1) -- (-2,-0.1);
        \draw (0.1,1) -- (-0.1,1);

        % Графік за точками: (0;4) (1;3) (2;2) (3;0) (4;-1)
        % Екстраполяція вліво до x=-2 для повноти картинки
        \draw[thick] plot [smooth, tension=0.6] coordinates {(-2, 4.5) (-1, 4.2) (0, 4) (1, 3) (2, 2) (3, 0) (4, -1)};

        \fill (3,0) circle (3pt); % Акцент на правильній відповіді (перетин з віссю x)
        \node[above right, fill=white, inner sep=1.5pt] at (1,3) {$y=f(x)$};
    \end{tikzpicture}\end{nmtfit}
    \end{flushright}
\end{minipage}

\nmtnobreak
\nopagebreak\nbvspace{0.3cm}
\answerTableTall{$(0; 3)$}{$(3; 4)$}{$(3; 0)$}{$(4; 0)$}{$(0; 4)$}
\par\penalty-20
\end{samepage}

\begin{samepage}
\noindent\zadnum \begin{minipage}[t]{0.55\textwidth}
На рисунку зображено графік функції $y=f(x)$, визначеної на проміжку $[-4; 5]$. Точка $(x_0; -2)$ належить графіку цієї функції. Визначте абсцису $x_0$ цієї точки. \nmtyear{2023}
\end{minipage}
\hfill
\begin{minipage}[t]{0.4\textwidth}
    \nopagebreak\nbvspace{-0.3cm}
    \begin{flushright}
    \begin{nmtfit}\begin{tikzpicture}[scale=0.5]
        \draw[step=1cm,gray!50,very thin] (-4.5,-3.5) grid (5.5,3.5);
        \draw[->, >=stealth, thick] (-4.5,0) -- (5.5,0) node[below] {$x$};
        \draw[->, >=stealth, thick] (0,-3.5) -- (0,3.5) node[left] {$y$};

        \node[below left] at (0,0) {$0$};
        \node[below] at (1,0) {$1$};
        \node[left] at (0,1) {$1$};
        \node[below] at (5,0) {$5$};
        \node[below] at (-4,0) {$-4$};

        \draw (1,0.1) -- (1,-0.1);
        \draw (5,0.1) -- (5,-0.1);
        \draw (-4,0.1) -- (-4,-0.1);
        \draw (0.1,1) -- (-0.1,1);

        % Графік за точками:
        % (-4,-3) (-3,-2) (-2,0) (-1,3) (0,2) (1,0.5) (2,-1) (3,0) (4,0.8) (5,1)
        \draw[thick] plot [smooth, tension=0.4] coordinates {(-4,-3) (-3,-2) (-2,0) (-1,3) (0,2) (1,0.5) (2,-1) (3,0) (4,0.8) (5,1)};

        \fill (-3,-2) circle (3pt); % Шукана точка
        \node[above right, fill=white, inner sep=1.5pt] at (-1,3) {$y=f(x)$};
    \end{tikzpicture}\end{nmtfit}
    \end{flushright}
\end{minipage}

\nmtnobreak
\nopagebreak\nbvspace{0.3cm}
\answerTableTall{$3$}{$-3$}{$0$}{$-2$}{$2$}
\par\penalty-20
\end{samepage}

\begin{samepage}
\noindent\zadnum Укажіть графік непарної функції. \nmtyear{2024}

\nmtnobreak
\nopagebreak\nbvspace{0.3cm}
\begin{center}
\begin{tabular}{|*{5}{>{\centering\arraybackslash}m{2.8cm}|}}
\hline
\textbf{А} & \textbf{Б} & \textbf{В} & \textbf{Г} & \textbf{Д} \\
\hline
\begin{nmtfit}\begin{tikzpicture}[scale=0.4]
    \draw[->, >=stealth] (-2,0) -- (2,0) node[below] {$x$};
    \draw[->, >=stealth] (0,-2) -- (0,2) node[left] {$y$};
    \node[below right] at (0,0) {$0$};
    \draw[thick] (-1.5,-1.5) -- (1.5,1.5);
\end{tikzpicture}\end{nmtfit} &
\begin{nmtfit}\begin{tikzpicture}[scale=0.4]
    \draw[->, >=stealth] (-2,0) -- (2,0) node[below] {$x$};
    \draw[->, >=stealth] (0,-2) -- (0,2) node[left] {$y$};
    \node[below left] at (0,0) {$0$};
    \draw[thick] (-3,1.5) -- (0,0) -- (1.9,1.5);
\end{tikzpicture}\end{nmtfit} &
\begin{nmtfit}\begin{tikzpicture}[scale=0.4]
    \draw[->, >=stealth] (-2,0) -- (2,0) node[below] {$x$};
    \draw[->, >=stealth] (0,-2) -- (0,2) node[left] {$y$};
    \node[below left] at (0,0) {$0$};
    % Парабола або схоже на модуль, симетричне відносно OY
    \draw[thick] (-1.5,1.5) -- (0,0) -- (1.5,1.5); 
\end{tikzpicture}\end{nmtfit} &
\begin{nmtfit}\begin{tikzpicture}[scale=0.4]
    \draw[->, >=stealth] (-2,0) -- (2,0) node[below] {$x$};
    \draw[->, >=stealth] (0,-2) -- (0,2) node[left] {$y$};
    \node[below right] at (0,0) {$0$};
    \draw[thick] (-2,1.5) -- (-1,0) -- (1,1);
\end{tikzpicture}\end{nmtfit} &
\begin{nmtfit}\begin{tikzpicture}[scale=0.4]
    \draw[->, >=stealth] (-2,0) -- (2,0) node[below] {$x$};
    \draw[->, >=stealth] (0,-2) -- (0,2) node[left] {$y$};
    \node[below left] at (0,0) {$0$};
    \draw[thick] (-1,-2) -- (1.5,1);
\end{tikzpicture}\end{nmtfit} \\
\hline
\end{tabular}
\end{center}
\par\penalty-20\vspace{0.25cm}
\end{samepage}

\begin{samepage}
\noindent\zadnum \begin{minipage}[t]{0.55\textwidth}
На рисунку зображено графік функції $y=f(x)$, визначеної на проміжку $[-4; 6]$. Укажіть різницю між найбільшим і найменшим значенням функції $f(x)$ на цьому проміжку. \nmtyear{2024}
\end{minipage}
\hfill
\begin{minipage}[t]{0.4\textwidth}
    \nopagebreak\nbvspace{-0.3cm}
    \begin{flushright}
    \begin{nmtfit}\begin{tikzpicture}[scale=0.5]
        % Розширив сітку: y від -2.5 до 6.5
        \draw[step=1cm,gray!50,very thin] (-4.5,-2.5) grid (6.5,6.5);
        \draw[->, >=stealth, thick] (-4.5,0) -- (6.5,0) node[below] {$x$};
        \draw[->, >=stealth, thick] (0,-2.5) -- (0,6.5) node[left] {$y$};

        \node[below left] at (0,0) {$0$};
        \node[below] at (1,0) {$1$};
        \node[left] at (0,1) {$1$};
        \node[below] at (6,0) {$6$};
        \node[below] at (-4,0) {$-4$};

        % Графік: (-4, 6) -> (-3, 0) -> (-2, -2) [мін] -> (1, 1) -> (4, 3) -> (6, 1)
        \draw[thick] plot [smooth, tension=0.6] coordinates {(-4, 6) (-3, 0) (-2, -2) (1, 1) (4, 3) (6, 1)};

        \fill (-4,6) circle (3pt);
        \fill (6,1) circle (3pt);
        \node[above right, fill=white, inner sep=1.5pt] at (4,3) {$y=f(x)$};
    \end{tikzpicture}\end{nmtfit}
    \end{flushright}
\end{minipage}

\nmtnobreak
\nopagebreak\nbvspace{0.3cm}
\answerTableTall{$6$}{$7$}{$5$}{$3$}{$10$}
\par\penalty-20
\end{samepage}

\begin{samepage}
\noindent\zadnum \begin{minipage}[t]{0.55\textwidth}
На рисунку зображено графік функції $y=f(x)$, визначеної на проміжку $[-5; 5]$. Укажіть поміж наведених координати точки, що належить цьому графіку. \nmtyear{2024}
\end{minipage}
\hfill
\begin{minipage}[t]{0.4\textwidth}
    \nopagebreak\nbvspace{-0.3cm}
    \begin{flushright}
    \begin{nmtfit}\begin{tikzpicture}[scale=0.45]
        \draw[step=1cm,gray!50,very thin] (-5.5,-2.5) grid (5.5,4.5);
        \draw[->, >=stealth, thick] (-5.5,0) -- (5.5,0) node[below] {$x$};
        \draw[->, >=stealth, thick] (0,-2.5) -- (0,4.5) node[left] {$y$};

        \node[below left] at (0,0) {$0$};
        \node[below] at (1,0) {$1$};
        \node[left] at (0,1) {$1$};
        \node[below] at (5,0) {$5$};
        \node[below] at (-5,0) {$-5$};

        % Графік через: (-5, 4), (-3, 2), (0, 0), (3, -2), (5, 1)
        \draw[thick] plot [smooth, tension=0.6] coordinates {(-5, 4) (-3, 2) (0, 0) (3, -2) (5, 1)};

        \fill (-5,4) circle (3pt);
        \fill (5,1) circle (3pt);
        \node[right, fill=white, inner sep=1.5pt] at (-3, 2.5) {$y=f(x)$};
    \end{tikzpicture}\end{nmtfit}
    \end{flushright}
\end{minipage}

\nmtnobreak
\nopagebreak\nbvspace{0.3cm}
\answerTableTall{$(2; -2)$}{$(-3; 2)$}{$(-2; 3)$}{$(-4; 3)$}{$(-1; 4)$}
\par\penalty-20
\end{samepage}

\begin{samepage}
\noindent\zadnum \begin{minipage}[t]{0.55\textwidth}
На рисунку зображено графік функції $y=f(x)$, визначеної на проміжку $[-5; 4]$. Скільки точок перетину з осями координат має ця функція на заданому проміжку? \nmtyear{2024}
\end{minipage}
\hfill
\begin{minipage}[t]{0.4\textwidth}
    \nopagebreak\nbvspace{-0.3cm}
    \begin{flushright}
    \begin{nmtfit}\begin{tikzpicture}[scale=0.5]
        \draw[step=1cm,gray!50,very thin] (-5.5,-2.5) grid (4.5,5.5);
        \draw[->, >=stealth, thick] (-5.5,0) -- (4.5,0) node[below] {$x$};
        \draw[->, >=stealth, thick] (0,-2.5) -- (0,5.5) node[left] {$y$};

        \node[below left] at (0,0) {$0$};
        \node[below] at (1,0) {$1$};
        \node[left] at (0,1) {$1$};
        \node[below] at (4,0) {$4$};
        \node[below] at (-5,0) {$-5$};

        % Графік
        \draw[thick] plot [smooth, tension=0.7] coordinates {(-5, -2) (-2, 3) (1, -2) (4, 5)};

        \fill (-5,-2) circle (3pt);
        \fill (4,5) circle (3pt);
        \node[left, fill=white, inner sep=1.5pt] at (-2, 2.5) {$y=f(x)$};
    \end{tikzpicture}\end{nmtfit}
    \end{flushright}
\end{minipage}

\nmtnobreak
\nopagebreak\nbvspace{0.3cm}
\answerTableTall{$3$}{$6$}{$2$}{$5$}{$4$}
\par\penalty-20
\end{samepage}

\begin{samepage}
\noindent\zadnum \begin{minipage}[t]{0.55\textwidth}
На рисунку зображено графік функції $y=f(x)$, визначеної на проміжку $[-5; 5]$. Укажіть різницю між найбільшим і найменшим значенням функції $f(x)$ на проміжку $[0; 5]$. \nmtyear{2024}
\end{minipage}
\hfill
\begin{minipage}[t]{0.4\textwidth}
    \nopagebreak\nbvspace{-0.3cm}
    \begin{flushright}
    \begin{nmtfit}\begin{tikzpicture}[scale=0.5]
        % Сітка до 7.5
        \draw[step=1cm,gray!50,very thin] (-5.5,-2.5) grid (5.5,7.5);
        \draw[->, >=stealth, thick] (-5.5,0) -- (5.5,0) node[below] {$x$};
        \draw[->, >=stealth, thick] (0,-2.5) -- (0,7.5) node[left] {$y$};

        \node[below left] at (0,0) {$0$};
        \node[below] at (1,0) {$1$};
        \node[left] at (0,1) {$1$};
        \node[below] at (5,0) {$5$};
        \node[below] at (-5,0) {$-5$};

        % Графік для x>=0: (0, 3), (2, 1), (3, 2), (4, 4), (5, 7)
        % Додали ліву частину для цілісності вигляду
        \draw[thick] plot [smooth, tension=0.25] coordinates {(-5, 6) (-3, 4) (0, 3) (2, 1) (3, 2) (4, 4) (5, 7)};

        \fill (-5,6) circle (3pt);
        \fill (5,7) circle (3pt);
        \node[left, fill=white, inner sep=1.5pt] at (-1, 3.5) {$y=f(x)$};
    \end{tikzpicture}\end{nmtfit}
    \end{flushright}
\end{minipage}

\nmtnobreak
\nopagebreak\nbvspace{0.3cm}
\answerTable{$10$}{$7$}{$3$}{$8$}{$6$}
\par\penalty-20
\end{samepage}

\begin{samepage}
\noindent\zadnum \begin{minipage}[t]{0.55\textwidth}
На рисунку зображено графік функції $y=f(x)$, визначеної на проміжку $[-5; 5]$. Укажіть \textit{ординату} точки перетину графіка функції з віссю $y$. \nmtyear{2025}
\end{minipage}
\hfill
\begin{minipage}[t]{0.4\textwidth}
    \nopagebreak\nbvspace{-0.3cm}
    \begin{flushright}
    \begin{nmtfit}\begin{tikzpicture}[scale=0.5]
        % Сітка
        \draw[step=1cm,gray!50,very thin] (-5.5,-2.5) grid (5.5,6.5);
        % Осі
        \draw[->, >=stealth, thick] (-5.5,0) -- (5.5,0) node[below] {$x$};
        \draw[->, >=stealth, thick] (0,-2.5) -- (0,6.5) node[left] {$y$};

        % Підписи
        \node[below left] at (0,0) {$0$};
        \node[below] at (1,0) {$1$};
        \node[left] at (0,1) {$1$};
        \node[below] at (5,0) {$5$};
        \node[below] at (-5,0) {$-5$};

        % Графік (координати підібрані візуально під картинку)
        \draw[thick] plot [smooth, tension=0.6] coordinates {(-5, -2) (-2, 4) (0, 3) (1, 2) (5, 6)};

        % Точки на кінцях
        \fill (-5,-2) circle (3pt);
        \fill (5,6) circle (3pt);
        \node[right, fill=white, inner sep=1.5pt] at (2, 4.5) {$y=f(x)$};
    \end{tikzpicture}\end{nmtfit}
    \end{flushright}
\end{minipage}

\nmtnobreak
\nopagebreak\nbvspace{0.3cm}
\answerTableTall{$3$}{$1$}{$0$}{$-4$}{$6$}
\par\penalty-20
\end{samepage}

\begin{samepage}
\noindent\zadnum \begin{minipage}[t]{0.55\textwidth}
На рисунку зображено графік функції $y=f(x)$, визначеної на відрізку $[-5; 5]$. Позначте властивість, яку має ця функція. \nmtyear{2025}
\end{minipage}
\hfill
\begin{minipage}[t]{0.4\textwidth}
    \nopagebreak\nbvspace{-0.3cm}
    \begin{flushright}
    \begin{nmtfit}\begin{tikzpicture}[scale=0.5]
        % Сітка
        \draw[step=1cm,gray!50,very thin] (-5.5,-3.5) grid (5.5,6.5);
        % Осі
        \draw[->, >=stealth, thick] (-5.5,0) -- (5.5,0) node[below] {$x$};
        \draw[->, >=stealth, thick] (0,-3.5) -- (0,6.5) node[left] {$y$};

        % Підписи
        \node[below left] at (0,0) {$0$};
        \node[below] at (1,0) {$1$};
        \node[left] at (0,1) {$1$};
        \node[below] at (5,0) {$5$};
        \node[below] at (-5,0) {$-5$};

        % Графік
        \draw[thick] plot [smooth, tension=0.6] coordinates {(-5, -3) (-1.5, 3.5) (0, 2) (1.5, 1) (5, 6)};

        % Точки на кінцях
        \fill (-5,-3) circle (3pt);
        \fill (5,6) circle (3pt);
        \node[right, fill=white, inner sep=1.5pt] at (2, 4.5) {$y=f(x)$};
    \end{tikzpicture}\end{nmtfit}
    \end{flushright}
\end{minipage}

\nmtnobreak
\nopagebreak\nbvspace{0.2cm}
\begin{tabular}{ll}
    \textbf{А} & набуває лише додатних значень \\[0.3cm]
    \textbf{Б} & парна \\[0.3cm]
    \textbf{В} & має дві точки екстремуму \\[0.3cm]
    \textbf{Г} & непарна \\[0.3cm]
    \textbf{Д} & зростає на всій області визначення \\
\end{tabular}
\par\penalty-20\vspace{0.25cm}
\end{samepage}

\begin{samepage}
\noindent\zadnum Укажіть парну функцію. \nmtyear{2025}

\nmtnobreak
\nopagebreak\nbvspace{0.3cm}
\answerTableTall{$y=2x^4$}{$y=\sin x + 1$}{$y=2x$}{$y=3^{2x}$}{$y=\dfrac{3}{x}$}
\par\penalty-20
\end{samepage}

\begin{samepage}
\noindent\zadnum \begin{minipage}[t]{0.55\textwidth}
На рисунку зображено графік функції $y=f(x)$, визначеної на проміжку $[-2; 4]$. Укажіть значення аргумента $x$, за якого функція набуває додатних значень. \nmtyear{2025}
\end{minipage}
\hfill
\begin{minipage}[t]{0.4\textwidth}
    \nopagebreak\nbvspace{-0.3cm}
    \begin{flushright}
    \begin{nmtfit}\begin{tikzpicture}[scale=0.5]
        % Сітка
        \draw[step=1cm,gray!50,very thin] (-2.5,-2.5) grid (4.5,4.5);
        % Осі
        \draw[->, >=stealth, thick] (-2.5,0) -- (4.5,0) node[below] {$x$};
        \draw[->, >=stealth, thick] (0,-2.5) -- (0,4.5) node[left] {$y$};

        % Підписи
        \node[below left] at (0,0) {$0$};
        \node[below] at (1,0) {$1$};
        \node[left] at (0,1) {$1$};
        \node[below] at (4,0) {$4$};
        \node[below] at (-2,0) {$-2$};

        % Графік
        \draw[thick] plot [smooth, tension=0.6] coordinates {(-2, 1) (-1, 3.5) (1, 0) (3, -1.5) (4, -2)};

        % Точки на кінцях
        \fill (-2,1) circle (3pt);
        \fill (4,-2) circle (3pt);
        \node[right, fill=white, inner sep=1.5pt] at (2, 2.5) {$y=f(x)$};
    \end{tikzpicture}\end{nmtfit}
    \end{flushright}
\end{minipage}

\nmtnobreak
\nopagebreak\nbvspace{0.2cm}
\answerTableTall{$4$}{$2$}{$0$}{$3$}{$1$}
\par\penalty-20
\end{samepage}

\begin{samepage}
\noindent\zadnum \begin{minipage}[t]{0.55\textwidth}
На рисунку зображено графік функції $y=f(x)$, визначеної на відрізку $[-5; 5]$. На якому з наведених проміжків функція має нуль? \nmtyear{2025}
\end{minipage}
\hfill
\begin{minipage}[t]{0.4\textwidth}
    \nopagebreak\nbvspace{-0.3cm}
    \begin{flushright}
    \begin{nmtfit}\begin{tikzpicture}[scale=0.5]
        % Сітка
        \draw[step=1cm,gray!50,very thin] (-5.5,-2.5) grid (5.5,5.5);
        % Осі
        \draw[->, >=stealth, thick] (-5.5,0) -- (5.5,0) node[below] {$x$};
        \draw[->, >=stealth, thick] (0,-2.5) -- (0,5.5) node[left] {$y$};

        % Підписи
        \node[below left] at (0,0) {$0$};
        \node[below] at (1,0) {$1$};
        \node[left] at (0,1) {$1$};
        \node[below] at (5,0) {$5$};
        \node[below] at (-5,0) {$-5$};

        % Графік
        \draw[thick] plot [smooth, tension=0.6] coordinates {(-5, -2) (-4, 0) (-2, 4) (0, 2) (1, 1) (5, 5)};

        % Точки на кінцях
        \fill (-5,-2) circle (3pt);
        \fill (5,5) circle (3pt);
        \node[right, fill=white, inner sep=1.5pt] at (2, 3.5) {$y=f(x)$};
    \end{tikzpicture}\end{nmtfit}
    \end{flushright}
\end{minipage}

\nmtnobreak
\nopagebreak\nbvspace{0.2cm}
\begin{tabular}{ll}
    \textbf{А} & $[-5; -3)$ \\[0.3cm]
    \textbf{Б} & $[-3; -1)$ \\[0.3cm]
    \textbf{В} & $[-1; 1)$ \\[0.3cm]
    \textbf{Г} & $[1; 3)$ \\[0.3cm]
    \textbf{Д} & $[3; 5]$ \\
\end{tabular}
\par\penalty-20\vspace{0.25cm}
\end{samepage}

\begin{samepage}
% НМТ 2026, сесія 2.06, завдання №8
\zadtask{Задано функцію $f(x) = x^2$. Укажіть правильну нерівність. \nmtyear{2026}}
\answerTable{$f(4) < 0$}{$f(-4) < f(4)$}{$f(-1) > f(2)$}{$f(-1) < 0$}{$f(3) < f(4)$}
% Відповідь: Д
\ifshowsolutions\solution{Для $f(x)=x^2$ маємо $f(3)=9$, $f(4)=16$, тож $f(3)<f(4)$ -- правильно. Відповідь: \textbf{Д}.}\fi
\par\penalty-20
\end{samepage}

\begin{samepage}
% НМТ 2026, сесія 3.06, завдання №8 --- також у темі: Графіки елементарних функцій, їх побудова та перетворення
\noindent
\begin{minipage}[c]{0.56\textwidth}
\zadnum У прямокутній системі координат зображено фрагмент графіка \textit{непарної} функції $y = f(x)$ та точки $K$, $L$, $M$, $N$, $P$. Через яку з цих точок проходить графік функції $y = f(x)$? \nmtyear{2026}
\end{minipage}%
\hfill
\begin{minipage}[c]{0.40\textwidth}
\centering
\begin{nmtfit}\begin{tikzpicture}[scale=0.5, thick]
    \draw[gray!35, very thin] (-3.3,-3.3) grid (3.3,3.3);
    \draw[->] (-3.5,0) -- (3.5,0) node[right, font=\scriptsize] {$x$};
    \draw[->] (0,-3.5) -- (0,3.5) node[above, font=\scriptsize] {$y$};
    \node[below left, font=\tiny] at (0,0) {$0$};
    \node[font=\tiny] at (1,-0.35) {$1$};
    \node[font=\tiny] at (-0.35,1) {$1$};
    % фрагмент у I чверті: відрізок від (1,1) до (2,3) приблизно
    \draw[mainGreen, line width=1.2pt] (1,1) -- (2,3);
    % точки
    \fill (-3,2) circle (2.4pt); \node[left, font=\tiny] at (-3,2) {$L$};
    \fill (-2.6,1.4) circle (2.4pt); \node[left, font=\tiny] at (-2.6,1.4) {$K$};
    \fill (-2,-3) circle (2.4pt); \node[left, font=\tiny] at (-2,-3) {$M$};
    \fill (2,-2.4) circle (2.4pt); \node[right, font=\tiny] at (2,-2.4) {$N$};
    \fill (-3,-3.1) circle (2.4pt); \node[left, font=\tiny] at (-3,-3.1) {$P$};
\end{tikzpicture}\end{nmtfit}
\end{minipage}
\par\nbvspace{0.2cm}
\answerTable{$K$}{$L$}{$M$}{$N$}{$P$}
% Відповідь: В
\ifshowsolutions\solution{Графік непарної функції симетричний відносно початку координат. Точка, симетрична до точки $(2;\,3)$ цього фрагмента,~--- це $(-2;\,-3)$, тобто точка $M$. Відповідь: \textbf{В} ($M$).}\fi
\par\penalty-20
\end{samepage}

\begin{samepage}
% НМТ 2026, сесія 4.06, завдання №8
\noindent
\begin{minipage}[c]{0.52\textwidth}
\zadnum На рисунку зображено графік функції $y = f(x)$, визначеної на проміжку $[-4;\, 4]$. Укажіть точку мінімуму цієї функції. \nmtyear{2026}
\end{minipage}%
\hfill
\begin{minipage}[c]{0.46\textwidth}
\centering
\begin{nmtfit}\begin{tikzpicture}[scale=0.55, thick]
    % Сітка
    \draw[gray!40, thin] (-4.5,-4.5) grid (6.5,6.5);

    % Осі
    \draw[->, thick] (-4.5,0) -- (4.8,0) node[below] {$x$};
    \draw[->, thick] (0,-2.5) -- (0,4.8) node[left] {$y$};

    % Підписи осей
    \node[below left] at (0,0) {$0$};
    \node[below] at (1,0) {$1$};
    \node[below] at (-4,0) {$-4$};
    \node[below] at (4,0) {$4$};
    \node[left] at (0,1) {$1$};

    % Графік функції (гладка крива по ключових точках)
    \draw[mainGreen!80!black, very thick] plot[smooth, tension=0.45] coordinates {
        (-4,-2) (-2.3,5) (0,3) (1.5,1) (3,-1) (4,0)
    };

    % Точки на кінцях
    \fill[mainGreen!80!black] (-4,-2) circle (3pt);
    \fill[mainGreen!80!black] (4,0) circle (3pt);

    % Підпис графіка у білому прямокутнику
    \node[fill=white, inner sep=1pt] at (4, 3.8) {$y = f(x)$};
\end{tikzpicture}\end{nmtfit}
\end{minipage}
\par\nbvspace{0.4cm}
\answerTable{$-2$}{$1$}{$3$}{$2$}{$5$}
% Відповідь: В
\ifshowsolutions\solution{Точка мінімуму~--- абсциса, у якій функція переходить від спадання до зростання (локальна «впадина»). На графіку це відбувається при $x=3$ (там $y=-1$). Відповідь: \textbf{В}.}\fi
\par\penalty-20
\end{samepage}

\begin{samepage}
% НМТ 2026, сесія 5.06, завдання №8 --- основна тема: Графіки елементарних функцій, їх побудова та перетворення
\noindent
\begin{minipage}[c]{0.56\textwidth}
\zadnum На рисунку зображено графік функції $y = f(x)$, визначеної на проміжку $[-5;\,5]$. Визначте точку екстремуму функції $y = f(x - 3)$. \nmtyear{2026}
\end{minipage}%
\hfill
\begin{minipage}[c]{0.40\textwidth}
\centering
\begin{nmtfit}\begin{tikzpicture}[scale=0.45, thick]
    \draw[gray!30, very thin] (-5.3,-2.3) grid (5.3,4.3);
    \draw[->] (-5.5,0) -- (5.6,0) node[right, font=\scriptsize] {$x$};
    \draw[->] (0,-2.5) -- (0,4.5) node[above, font=\scriptsize] {$y$};
    \node[below left, font=\tiny] at (0,0) {$0$};
    \node[font=\tiny] at (1,-0.4) {$1$}; \node[font=\tiny] at (-5,-0.4) {$-5$}; \node[font=\tiny] at (5,-0.4) {$5$};
    \draw[mainGreen, line width=1.2pt]
        (-5,-2) .. controls (-3,0) and (0,3.5) .. (2,4)
        .. controls (3.5,4) and (4.5,2) .. (5,1.5);
    \node[font=\tiny, text=mainGreen, fill=white, inner sep=1.5pt] at (3.5,4) {$y=f(x)$};
\end{tikzpicture}\end{nmtfit}
\end{minipage}
\par\nbvspace{0.2cm}
\answerTable{$5$}{$1$}{$-1$}{$2$}{$-5$}
% Відповідь: А
\ifshowsolutions\solution{Графік $y=f(x)$ має екстремум (максимум) у точці $x=2$. Заміна $x\to x-3$ зсуває графік праворуч на $3$, тож точка екстремуму $x=2+3=5$. Відповідь: \textbf{А}.}\fi
\par\penalty-20
\end{samepage}

\begin{samepage}
% НМТ 2026, сесія 9.06, завдання №8
\noindent
\begin{minipage}[c]{0.56\textwidth}
\zadnum На рисунку зображено графік функції $y=f(x)$, визначеної на проміжку $[-4;\, 4]$. Укажіть проміжок, на якому функція набуває лише від'ємних значень. \nmtyear{2026}
\end{minipage}%
\hfill
\begin{minipage}[c]{0.42\textwidth}
\centering
\begin{nmtfit}\begin{tikzpicture}[scale=0.6, thick]
    \draw[step=1cm, gray!40, thin] (-4,-2) grid (5,3);
    \draw[->] (-4.5,0) -- (5.5,0) node[right] {$x$};
    \draw[->] (0,-2.5) -- (0,3.5) node[above] {$y$};
    \node[below left] at (0,0) {$0$};
    \node[below right] at (1,0) {$1$};
    \node[above left] at (0,1) {$1$};
    \node[below] at (-4,0) {$-4$};
    \node[below] at (4,0) {$4$};
    \draw[magenta!80!black, thick, smooth] plot coordinates {(-4,-1) (-2,1.5) (0,1) (1,0) (2,-1) (3,0) (4,1)};
    \fill (-4,-1) circle (3pt);
    \fill (4,1) circle (3pt);
    \node[above right, fill=white, inner sep=1.5pt] at (1,1.5) {$y=f(x)$};
\end{tikzpicture}\end{nmtfit}
\end{minipage}
\par\nbvspace{0.2cm}
\answerTable{$(-4; -1)$}{$(0; 1)$}{$(1; 2)$}{$(1; 4)$}{$(3; 4)$}
\par\penalty-20
\end{samepage}

\begin{samepage}
% НМТ 2026, сесія 10.06, завдання №8
\noindent
\begin{minipage}[c]{0.56\textwidth}
\zadnum На рисунку зображено графік функції $y=f(x)$, визначеної на проміжку $(-\infty;\, +\infty)$. Скористайтеся рисунком і визначте нулі цієї функції. \nmtyear{2026}
\end{minipage}%
\hfill
\begin{minipage}[c]{0.42\textwidth}
\centering
\begin{nmtfit}\begin{tikzpicture}[scale=0.45, thick]
    \draw[step=1cm, gray!40, thin] (-1,-2) grid (6,9);
    \draw[->] (-1.5,0) -- (6.5,0) node[right] {$x$};
    \draw[->] (0,-2.5) -- (0,9.5) node[above] {$y$};
    \node[below left, font=\small] at (0,0) {$0$};
    \node[below right, font=\small] at (1,0) {$1$};
    \node[above left, font=\small] at (0,1) {$1$};
    \node[left, font=\small] at (0,8) {$8$};
    \node[below, font=\small] at (4,0) {$4$};
    \node[below, font=\small] at (2,0) {$2$};
    \draw[domain=-0.35:6.35, smooth, samples=80, variable=\x, black, thick] plot ({\x}, {\x*\x-6*\x+8});
    \node[above right, font=\small, fill=white, inner sep=1.5pt] at (4.5,5) {$y=f(x)$};
\end{tikzpicture}\end{nmtfit}
\end{minipage}
\par\nbvspace{0.2cm}
\answerTable{$0;\ 8$}{$0$}{$5$}{$2;\ 4$}{$8$}
\par\penalty-20
\end{samepage}

\begin{samepage}
% НМТ 2026, сесія 15.06, завдання №8
\noindent
\begin{minipage}[c]{0.58\textwidth}
\zadnum На рисунку зображено графік функції $y = f(x)$, визначеної на проміжку $[-3;\, 3]$. На якому з наведених проміжків ця функція спадає? \nmtyear{2026}
\end{minipage}%
\hfill
\begin{minipage}[c]{0.40\textwidth}
\centering
\begin{nmtfit}\begin{tikzpicture}[scale=0.5, thick]
    \draw[step=1cm, gray!40, thin] (-3,-2) grid (4,5);
    \draw[->] (-3.3,0) -- (4.3,0) node[right, font=\scriptsize] {$x$};
    \draw[->] (0,-2.3) -- (0,5.3) node[above, font=\scriptsize] {$y$};
    \node[below left, font=\scriptsize] at (0,0) {$0$};
    \node[below, font=\scriptsize] at (-3,0) {$-3$};
    \node[below, font=\scriptsize] at (1,0) {$1$};
    \node[below, font=\scriptsize] at (3,0) {$3$};
    \node[left, font=\scriptsize] at (0,1) {$1$};
    \draw[mainGreen, very thick, smooth] plot coordinates {(-3,-2) (-1,0.5)(0,3.2) (1,4) (2,3.2) (3,1)};
    \node[right, font=\scriptsize, fill=white, inner sep=1.5pt] at (1.5,4.6) {$y=f(x)$};
\end{tikzpicture}\end{nmtfit}
\end{minipage}
\par\nbvspace{0.2cm}
\answerTable{$[-3; 1]$}{$[-3; 4]$}{$[2; 4]$}{$[1; 3]$}{$[-3; 3]$}
% Відповідь: Г
\par\penalty-20
\end{samepage}

\begin{samepage}
% НМТ 2026, сесія 16.06, завдання №8
\noindent
\begin{minipage}[c]{0.50\textwidth}
\zadnum На рисунку зображено графік функції $y=f(x)$, визначеної на відрізку $[-5;\, 5]$. На якому проміжку ця функція спадає? \nmtyear{2026}\par
\nopagebreak\nbvspace{0.6cm}

\textbf{А}\quad $[-2;\, 2]$\\
\textbf{Б}\quad $[-5;\, -2]$\\
\textbf{В}\quad $[1;\, 6]$\\
\textbf{Г}\quad $[-3;\, 2]$\\
\textbf{Д}\quad $[2;\, 6]$
\end{minipage}%
\hfill
\begin{minipage}[c]{0.48\textwidth}
\centering
\begin{nmtfit}\begin{tikzpicture}[scale=0.6, thick]
    \draw[step=1cm, gray!50, thin] (-6,-4) grid (6,4);
    \draw[->, >=latex, very thick] (-6,0) -- (6.5,0) node[below] {$x$};
    \draw[->, >=latex, very thick] (0,-4) -- (0,4.5) node[left] {$y$};

    \node[below left] at (0,0) {$0$};
    \node[below] at (-5,0) {$-5$};
    \node[below] at (5,0) {$5$};
    \node[below] at (1,0) {$1$};
    \node[left] at (0,1) {$1$};

    % Графік функції
    \draw[line width=1.2pt, smooth, tension=0.7] plot coordinates {(-5,-3) (-4,0) (-2,3) (0,1) (2,-1) (4,1) (5,3)};

    \fill (-5,-3) circle (3pt);
    \fill (5,3) circle (3pt);

    \node[above right, fill=white, inner sep=1pt] at (1.5, 1.5) {$y=f(x)$};
\end{tikzpicture}\end{nmtfit}
\end{minipage}
% Відповідь: А
\par\penalty-20
\end{samepage}

\begin{samepage}
% НМТ 2026, сесія 17.06, завдання №8
\noindent
\begin{minipage}[c]{0.58\textwidth}
\zadnum На рисунку зображено графік квадратичної функції. Укажіть точку локального екстремуму цієї функції. \nmtyear{2026}
\end{minipage}%
\hfill
\begin{minipage}[c]{0.40\textwidth}
\centering
\begin{nmtfit}\begin{tikzpicture}[scale=0.5, thick]
    \draw[step=1cm, gray!40, thin] (-2,-3) grid (4,4);
    \draw[->] (-2.3,0) -- (4.3,0) node[right, font=\scriptsize] {$x$};
    \draw[->] (0,-3.3) -- (0,4.3) node[above, font=\scriptsize] {$y$};
    \node[below left, font=\scriptsize] at (0,0) {$0$};
    \node[below, font=\scriptsize] at (1,0) {$1$};
    \node[left, font=\scriptsize] at (0,1) {$1$};
    % парабола гілками вниз, вершина (1,3): y = -(x-1)^2+3 -> прибл
    \draw[mainGreen, very thick, domain=-1.05:3.05, samples=60] plot (\x, {-1*(\x-1)*(\x-1)+3});
\end{tikzpicture}\end{nmtfit}
\end{minipage}
\par\nbvspace{0.2cm}
\answerTable{$x_0 = 4$}{$x_0 = 0$}{$x_0 = 1$}{$x_0 = 3$}{$x_0 = -1$}
% Відповідь: В
\par\penalty-20
\end{samepage}

\begin{samepage}
% НМТ 2026, сесія 19.06, завдання №8
\zadtask{Графік парної функції $y = f(x)$ проходить через точку $(-3;\, 9)$. Визначте $f(3)$. \nmtyear{2026}}
\answerTable{$9$}{$6$}{$3$}{$-9$}{$-3$}
% Відповідь: А
\par\penalty-20
\end{samepage}

\typeTitle{Завдання на встановлення відповідності}

\begin{samepage}
\noindent\zadnum До кожного початку речення \mbox{(1--3)} доберіть його закінчення \mbox{(А--Д)} так, щоб утворилося правильне твердження. \nmtyear{2023}

\nmtnobreak
\nopagebreak\nbvspace{0.3cm}

\nmtnobreak
\noindent
\matchingLayout{
\matchHead{Початок речення}
\matchItem{1}{Функція $y = \log_{0{,}5} x$}
\matchItem{2}{Функція $y = \sin x$}
\matchItem{3}{Функція $y = \dfrac{1}{2x-2}$}

\nmtnobreak
    % Зменшена табличка
}{
\matchHead{Закінчення речення}
    \matchItem{А}{не визначена при $x=1$.}
\matchItem{Б}{набуває від’ємного значення при $x=2$.}
\matchItem{В}{є непарною.}
\matchItem{Г}{має лише одну точку локального екстремуму.}
\matchItem{Д}{зростає на проміжку $(0; +\infty)$.}
}{
\matchingGrid
}
\par\penalty-20\vspace{0.25cm}
\end{samepage}

\begin{samepage}
\noindent\zadnum Установіть відповідність між функцією \mbox{(1--3)} та її властивістю \mbox{(А--Д)}. \nmtyear{2023}

\nmtnobreak
\nopagebreak\nbvspace{0.3cm}

\nmtnobreak
\noindent
\matchingLayout{
\matchHead{Функція}
\matchItem{1}{$y = 7x + 4$}
\matchItem{2}{$y = -\dfrac{7}{x}$}
\matchItem{3}{$y = \log_{0{,}5} (x-4)$}

\nmtnobreak
    % Зменшена табличка
}{
\matchHead{Властивість}
    \matchItem{А}{є спадною на всій області визначення}
\matchItem{Б}{графік функції перетинає вісь $y$ в точці з ординатою $4$}
\matchItem{В}{є непарною}
\matchItem{Г}{є парною}
\matchItem{Д}{областю визначення є проміжок $(0; +\infty)$}
}{
\matchingGrid
}
\par\penalty-20\vspace{0.25cm}
\end{samepage}

\begin{samepage}
\noindent\zadnum До кожного початку речення \mbox{(1--3)} доберіть його закінчення \mbox{(А--Д)} так, щоб утворилося правильне твердження. \nmtyear{2023}

\nmtnobreak
\nopagebreak\nbvspace{0.3cm}

\nmtnobreak
\noindent
\matchingLayout{
\matchHead{Початок речення}
\matchItem{1}{Функція $y = \log_2 x$}
\matchItem{2}{Функція $y = x^2 - 4x - 4$}
\matchItem{3}{Функція $y = \dfrac{1}{x}$}

\nmtnobreak
    % Зменшена табличка
}{
\matchHead{Закінчення речення}
    \matchItem{А}{не визначена при $x = -2$.}
\matchItem{Б}{набуває від’ємного значення при $x = 2$.}
\matchItem{В}{є непарною.}
\matchItem{Г}{спадає на проміжку $(-\infty; 4]$.}
\matchItem{Д}{зростає на проміжку $(-\infty; +\infty)$.}
}{
\matchingGrid
}
\par\penalty-20\vspace{0.25cm}
\end{samepage}

\begin{samepage}
\noindent\zadnum Установіть відповідність між функцією \mbox{(1--3)} та її властивістю \mbox{(А--Д)}. \nmtyear{2023}

\nmtnobreak
\nopagebreak\nbvspace{0.3cm}

\nmtnobreak
\noindent
\matchingLayout{
\matchHead{Функція}
\matchItem{1}{$y = x^2 - 1$}
\matchItem{2}{$y = 3^x$}
\matchItem{3}{$y = x^3$}

\nmtnobreak
    % Зменшена табличка
}{
\matchHead{Властивість}
    \matchItem{А}{проходить через початок координат}
\matchItem{Б}{не містить точок $(x_0; y_0)$ з від'ємною ординатою $y_0$}
\matchItem{В}{не має спільних точок з віссю $y$}
\matchItem{Г}{двічі перетинає графік прямої $y = 3$}
\matchItem{Д}{графік функції знаходиться лише в I координатній чверті}
}{
\matchingGrid
}
\par\penalty-20\vspace{0.25cm}
\end{samepage}

\begin{samepage}
\noindent\zadnum Установіть відповідність між твердженням \mbox{(1--3)} та функцією \mbox{(А--Д)}, для якої це твердження є правильним. \nmtyear{2023}

\nmtnobreak
\nopagebreak\nbvspace{0.3cm}

\nmtnobreak
\noindent
\matchingLayout{
% Трохи ширше для довгих тверджень
\matchHead{Твердження}
    \matchItem{1}{областю визначення функції є проміжок $[-1; +\infty)$}
\matchItem{2}{графік функції проходить через точку $(0; 1)$}
\matchItem{3}{функція має точку локального екстремуму на проміжку $[1; 3]$}

\nmtnobreak
    % Зменшена табличка
}{
\matchHead{Функція}
    \matchItem{А}{$y = -\cos x$}
\matchItem{Б}{$y = 4^x$}
\matchItem{В}{$y = 2\sqrt{x+1}$}
\matchItem{Г}{$y = x^2 - 4x + 3$}
\matchItem{Д}{$y = x$}
}{
\matchingGrid
}
\par\penalty-20\vspace{0.25cm}
\end{samepage}

\begin{samepage}
\noindent\zadnum Установіть відповідність між твердженням \mbox{(1--3)} та функцією \mbox{(А--Д)}, для якої це твердження є правильним. \nmtyear{2023}

\nmtnobreak
\nopagebreak\nbvspace{0.3cm}

\nmtnobreak
\noindent
\matchingLayout{
\matchHead{Твердження}
    \matchItem{1}{областю значень функції є проміжок $[0; +\infty)$}
\matchItem{2}{графік функції симетричний відносно осі $y$}
\matchItem{3}{найменше значення на відрізку $[1; 4]$ функція набуває в точці $x = 4$}

\nmtnobreak
    % Зменшена табличка
}{
\matchHead{Функція}
    \matchItem{А}{$y = x^2 + 4$}
\matchItem{Б}{$y = x$}
\matchItem{В}{$y = \sqrt{x}$}
\matchItem{Г}{$y = \log_{0{,}5} x$}
\matchItem{Д}{$y = -\dfrac{1}{x}$}
}{
\matchingGrid
}
\par\penalty-20\vspace{0.25cm}
\end{samepage}

\begin{samepage}
\noindent\zadnum Установіть відповідність між функцією \mbox{(1--3)} та її властивістю \mbox{(А--Д)}. \nmtyear{2023}

\nmtnobreak
\nopagebreak\nbvspace{0.3cm}

\nmtnobreak
\noindent
\matchingLayout{
\matchHead{Функція}
\matchItem{1}{$y = \sqrt{x+1}$}
\matchItem{2}{$y = 4-x^2$}
\matchItem{3}{$y = 3^{-x}$}

\nmtnobreak
    % Зменшена табличка
}{
\matchHead{Властивість}
    \matchItem{А}{має точку локального максимуму}
\matchItem{Б}{має точку локального мінімуму}
\matchItem{В}{є непарною}
\matchItem{Г}{зростає на всій області визначення}
\matchItem{Д}{набуває лише додатних значень}
}{
\matchingGrid
}
\par\penalty-20\vspace{0.25cm}
\end{samepage}

\begin{samepage}
\noindent\zadnum Установіть відповідність між функцією \mbox{(1--3)} та її властивістю \mbox{(А--Д)}. \nmtyear{2023}

\nmtnobreak
\nopagebreak\nbvspace{0.3cm}

\nmtnobreak
\noindent
\matchingLayout{
\matchHead{Функція}
\matchItem{1}{$y = (x+2)^2$}
\matchItem{2}{$y = 2\sqrt{x}$}
\matchItem{3}{$y = 2^x$}

\nmtnobreak
    % Зменшена табличка
}{
\matchHead{Властивість}
    \matchItem{А}{є зростаючою на проміжку $(-\infty; +\infty)$}
\matchItem{Б}{графік функції має одну спільну точку із графіком функції $y = x-2$}
\matchItem{В}{графік функції має дві спільні точки із графіком функції $y = x-2$}
\matchItem{Г}{є спадною на проміжку $(-\infty; -2]$}
\matchItem{Д}{є спадною на проміжку $(-\infty; 2]$}
}{
\matchingGrid
}
\par\penalty-20\vspace{0.25cm}
\end{samepage}

\begin{samepage}
\noindent\zadnum Установіть відповідність між функцією \mbox{(1--3)} та її властивістю \mbox{(А--Д)}. \nmtyear{2024}

\nmtnobreak
\nopagebreak\nbvspace{0.3cm}

\nmtnobreak
\noindent
\matchingLayout{
\matchHead{Функція}
\matchItem{1}{$y = 4 - x^2$}
\matchItem{2}{$y = -x^3$}
\matchItem{3}{$y = 4^x$}

\nmtnobreak
    % Табличка
}{
\matchHead{Властивість}
    \matchItem{А}{є зростаючою на всій області визначення}
\matchItem{Б}{набуває від’ємного значення при $x = -2$}
\matchItem{В}{є парною}
\matchItem{Г}{має три спільні точки з графіком функції $y = -x$}
\matchItem{Д}{графік функції розташований лише в I та IV координатній чверті}
}{
\matchingGrid
}
\par\penalty-20\vspace{0.25cm}
\end{samepage}

\begin{samepage}
\noindent\zadnum На рисунку зображено графік функції $y=f(x)$, визначеної на проміжку $[-4; 4]$. Установіть відповідність між початком речення \mbox{(1--3)} та його закінченням \mbox{(А--Д)} так, щоб утворилося правильне твердження. \nmtyear{2024}

\nmtnobreak
\nopagebreak\nbvspace{0.3cm}

\nmtnobreak
% --- ВЕРХНІЙ БЛОК: Умови + Графік ---
\noindent
\begin{minipage}[t]{0.55\textwidth}
    \textit{Початок речення} \par\par\nbvspace{0.3cm}
    \textbf{1} \quad Найменше значення функції $y=f(x)$ \\[0.4cm]
    \textbf{2} \quad Точка екстремуму функції $y=f(x)-5$ \\[0.4cm]
    \textbf{3} \quad Нуль функції $y=f(x+2)$
\end{minipage}%
\hfill
\begin{minipage}[t]{0.40\textwidth}
    \nopagebreak\nbvspace{-0.3cm} % Підтягуємо графік трохи вгору
    \begin{flushright}
    \begin{nmtfit}\begin{tikzpicture}[scale=0.5]
        % Сітка та осі
        \draw[step=1cm,gray!50,very thin] (-4.5,-3.5) grid (4.5,3.5);
        \draw[->, >=stealth, thick] (-4.5,0) -- (4.5,0) node[below] {$x$};
        \draw[->, >=stealth, thick] (0,-3.5) -- (0,3.5) node[left] {$y$};

        % Підписи осей
        \node[below left] at (0,0) {$0$};
        \node[below] at (1,0) {$1$};
        \node[left] at (0,1) {$1$};
        \node[below] at (4,0) {$4$};
        \node[below] at (-4,0) {$-4$};

        % Графік: старт (-4, -3), через (-2, -1), нуль (0,0), макс (3, 3), кінець (4, 1)
        \draw[thick] plot [smooth, tension=0.7] coordinates {(-4, -3) (-2, -1) (0, 0) (3, 3) (4, 1)};

        % Точки на кінцях
        \fill (-4,-3) circle (3pt);
        \fill (4,1) circle (3pt);

        \node[left, fill=white, inner sep=1.5pt] at (-1, 1) {$y=f(x)$};
    \end{tikzpicture}\end{nmtfit}
    \end{flushright}
\end{minipage}

\nmtnobreak
\nopagebreak\nbvspace{0.2cm}

\nmtnobreak
% --- НИЖНІЙ БЛОК: Варіанти + Таблиця ---
\noindent
\begin{minipage}[t]{0.55\textwidth}
    \textit{Закінчення речення} \par\par\nbvspace{0.2cm}
    \begin{tabular}{ll}
    \textbf{А} & дорівнює $-3$. \\
    \textbf{Б} & дорівнює $-2$. \\
    \textbf{В} & дорівнює $0$. \\
    \textbf{Г} & дорівнює $2$. \\
    \textbf{Д} & дорівнює $3$. \\
    \end{tabular}
\end{minipage}%
\hfill
\begin{minipage}[t]{0.40\textwidth}
    \nopagebreak\nbvspace{0.5cm} % Вирівнювання таблички по висоті
    \begin{flushright}
    % Табличка
    \begingroup
    \setlength{\tabcolsep}{4pt}
    \renewcommand{\arraystretch}{1.2}
    \small
    \begin{tabular}{r|c|c|c|c|c|}
         \multicolumn{1}{c}{} & \multicolumn{1}{c}{\textbf{А}} & \multicolumn{1}{c}{\textbf{Б}} & \multicolumn{1}{c}{\textbf{В}} & \multicolumn{1}{c}{\textbf{Г}} & \multicolumn{1}{c}{\textbf{Д}} \\ \cline{2-6}
         \textbf{1} & & & & & \\ \cline{2-6}
         \textbf{2} & & & & & \\ \cline{2-6}
         \textbf{3} & & & & & \\ \cline{2-6}
    \end{tabular}
    \endgroup
    \end{flushright}
\end{minipage}
\par\penalty-20\vspace{0.25cm}
\end{samepage}

\begin{samepage}
\noindent\zadnum Установіть відповідність між твердженням \mbox{(1--3)} та функцією \mbox{(А--Д)}, для якої це твердження є правильним. \nmtyear{2024}

\nmtnobreak
\nopagebreak\nbvspace{0.3cm}

\nmtnobreak
\noindent
\matchingLayout{
\matchHead{Твердження}
    \matchItem{1}{функція має 2 нулі}
\matchItem{2}{на відрізку $[-1; 3]$ функція набуває від'ємних значень}
\matchItem{3}{найменше значення функції на відрізку $[-1; 3]$ дорівнює $0{,}5$}

\nmtnobreak
    % Табличка
}{
\matchHead{Функція}
    \matchItem{А}{$y = x^2 - 4$}
\matchItem{Б}{$y = \dfrac{1}{x-4}$}
\matchItem{В}{$y = 2^x$}
\matchItem{Г}{$y = 0{,}5^x$}
\matchItem{Д}{$y = \sqrt{x+1}$}
}{
\matchingGrid
}
\par\penalty-20\vspace{0.25cm}
\end{samepage}

\begin{samepage}
\noindent\zadnum На рисунку зображено графік функції $y=f(x)$, визначеної на проміжку $[-4; 5]$. Установіть відповідність між початком речення \mbox{(1--3)} та його закінченням \mbox{(А--Д)} так, щоб утворилося правильне твердження. \nmtyear{2024}

\nmtnobreak
\nopagebreak\nbvspace{0.3cm}

\nmtnobreak
\noindent
\matchingLayout{
\matchHead{Початок речення}
\matchItem{1}{Нуль функції належить проміжку}
\matchItem{2}{Точка максимуму функції належить проміжку}
\matchItem{3}{Абсциса точки перетину графіка функції з графіком функції $y = \log_{\frac{1}{3}} x$ належить проміжку}
}{
\matchHead{Закінчення речення}
    \matchItem{А}{$(-4; -2]$.}
\matchItem{Б}{$(-2; 0]$.}
\matchItem{В}{$(0; 1]$.}
\matchItem{Г}{$(1; 3]$.}
\matchItem{Д}{$(3; 5]$.}
}{
\matchingGrid
}
\par\penalty-20\vspace{0.25cm}
\end{samepage}

\begin{samepage}
\noindent\zadnum На рисунку зображено графік функції $y=f(x)$, визначеної на проміжку $[1; 9]$. Установіть відповідність між початком речення \mbox{(1--3)} та його закінченням \mbox{(А--Д)} так, щоб утворилося правильне твердження. \nmtyear{2024}

\nmtnobreak
\nopagebreak\nbvspace{0.3cm}

\nmtnobreak
\noindent
\matchingLayout{
\matchHead{Початок речення}
\matchItem{1}{Найбільше значення функції на проміжку $[1; 9]$}
\matchItem{2}{Найменше значення функції на проміжку $[1; 3]$}
\matchItem{3}{Найменше ціле значення $x$, за якого виконується нерівність $f(x) < 0$}
}{
\matchHead{Закінчення речення}
    \matchItem{А}{дорівнює $5$.}
\matchItem{Б}{дорівнює $6$.}
\matchItem{В}{дорівнює $7$.}
\matchItem{Г}{дорівнює $8$.}
\matchItem{Д}{дорівнює $9$.}
}{
\matchingGrid
}
\par\penalty-20\vspace{0.25cm}
\end{samepage}

\begin{samepage}
\noindent\zadnum Установіть відповідність між твердженням \mbox{(1--3)} та функцією \mbox{(А--Д)}, для якої це твердження є правильним. \nmtyear{2024}

\nmtnobreak
\nopagebreak\nbvspace{0.3cm}

\nmtnobreak
\noindent
\matchingLayout{
\matchHead{Твердження}
    \matchItem{1}{є непарною}
\matchItem{2}{зростає на відрізку $[1; 4]$}
\matchItem{3}{найбільше значення функції на відрізку $[1; 4]$ є від'ємним числом}
}{
\matchHead{Функція}
    \matchItem{А}{$y = \dfrac{1}{x+1}$}
\matchItem{Б}{$y = \sin x$}
\matchItem{В}{$y = x^2 - 1$}
\matchItem{Г}{$y = 0{,}5^x$}
\matchItem{Д}{$y = -\sqrt{x}$}
}{
\matchingGrid
}
\par\penalty-20\vspace{0.25cm}
\end{samepage}

\begin{samepage}
\noindent\zadnum Узгодьте функцію (1–3) із її властивістю (А – Д). \nmtyear{2025}

\nmtnobreak
\nopagebreak\nbvspace{0.3cm}

\nmtnobreak
\noindent
\matchingLayout{
\matchHead{Функція}
    % Тут формули короткі, тому p{5cm} вистачить
    \matchItem{1}{$y=-x^3$}
\matchItem{2}{$y=\sqrt{x}$}
\matchItem{3}{$y=\cos x$}
}{
\matchHead{Властивість функції}
    % Тут текст довгий, тому використовуємо p{6.2cm} для перенесення слів
    \matchItem{А}{є парною}
\matchItem{Б}{має лише одну спільну точку з колом $x^2+y^2=9$}
\matchItem{В}{немає спільних точок з прямою $x=0$}
\matchItem{Г}{графік функції розташований лише в другій координатній чверті}
\matchItem{Д}{набуває всіх значень з проміжку $(-\infty; +\infty)$}
}{
\matchingGrid
}
\par\penalty-20\vspace{0.25cm}
\end{samepage}

\begin{samepage}
\noindent\zadnum Узгодьте функцію (1–3) із її властивістю (А – Д). \nmtyear{2025}

\nmtnobreak
\nopagebreak\nbvspace{0.3cm}

\nmtnobreak
\noindent
\matchingLayout{
\matchHead{Функція}
    \matchItem{1}{$y=x^2$}
\matchItem{2}{$y=x^3+1$}
\matchItem{3}{$y=\log_2 x$}
}{
\matchHead{Властивість функції}
    \matchItem{А}{парна}
\matchItem{Б}{непарна}
\matchItem{В}{область визначення функції є проміжок $(0; +\infty)$}
\matchItem{Г}{похідна функції є додатною на проміжку $(-\infty; 0)$}
\matchItem{Д}{є спадною на всій області визначення}
}{
\matchingGrid
}
\par\penalty-20\vspace{0.25cm}
\end{samepage}

\begin{samepage}
\noindent\zadnum Узгодьте функцію (1–3) із її властивістю (А – Д). \nmtyear{2025}

\nmtnobreak
\nopagebreak\nbvspace{0.3cm}

\nmtnobreak
\noindent
\matchingLayout{
\matchHead{Функція}
    \matchItem{1}{$y=7x+4$}
\matchItem{2}{$y=-\dfrac{7}{x}$}
\matchItem{3}{$y=\log_{0{,}1}(x-4)$}
}{
\matchHead{Властивість функції}
    \matchItem{А}{парна}
\matchItem{Б}{непарна}
\matchItem{В}{область визначення функції є проміжок $(0; +\infty)$}
\matchItem{Г}{графік функції перетинає вісь $y$ в точці з ординатою $4$}
\matchItem{Д}{спадає на всій області визначення}
}{
\matchingGrid
}
\par\penalty-20\vspace{0.25cm}
\end{samepage}

\begin{samepage}
\noindent\zadnum Увідповідніть функцію (1–3) та її властивість (А – Д). \nmtyear{2025}

\nmtnobreak
\nopagebreak\nbvspace{0.3cm}

\nmtnobreak
\noindent
\matchingLayout{
\matchHead{Функція}
    \matchItem{1}{$f(x)=x^2+2$}
\matchItem{2}{$f(x)=\sin x$}
\matchItem{3}{$f(x)=x^3$}
}{
\matchHead{Властивість функції}
    \matchItem{А}{зростає на всій області визначення}
\matchItem{Б}{спадає на всій області визначення}
\matchItem{В}{парна}
\matchItem{Г}{має більше ніж два нулі}
\matchItem{Д}{графік функції симетричний відносно осі $x$}
}{
\matchingGrid
}
\par\penalty-20\vspace{0.25cm}
\end{samepage}

\begin{samepage}
\noindent\zadnum На рисунку зображено графік функції $y=f(x)$, визначеної на проміжку $[-4; 5]$. Узгодьте проміжок (1–3) з властивістю (А – Д), яку має функція на цьому проміжку. \nmtyear{2025}

\nmtnobreak
\nopagebreak\nbvspace{0.2cm}

\nmtnobreak
\noindent
\begin{minipage}[t]{0.55\textwidth}
    \nopagebreak\nbvspace{0pt}
\matchHead{Проміжок}
    \matchItem{1}{$[-4; -2]$}
\matchItem{2}{$[-2; 2]$}
\matchItem{3}{$[2; 5]$}

    \nopagebreak\nbvspace{0.3cm}
\matchHead{Властивість функції}
    \matchItem{А}{функція зростає}
\matchItem{Б}{функція має точку мінімуму}
\matchItem{В}{функція спадає}
\matchItem{Г}{графік функції є фрагментом графіка функції $y=1-x^2$}
\matchItem{Д}{функція набуває лише додатних значень}

    \nopagebreak\nbvspace{0.3cm}

    \matchingGrid
\end{minipage}%
\hfill
\begin{minipage}[t]{0.40\textwidth}
    \nopagebreak\nbvspace{0pt}
    \begin{flushright}
    \begin{nmtfit}\begin{tikzpicture}[scale=0.5]
        \draw[step=1cm,gray!50,very thin] (-6.5,-7.5) grid (5.5,4.5);
        \draw[->, >=stealth, thick] (-6.5,0) -- (5.5,0) node[below] {$x$};
        \draw[->, >=stealth, thick] (0,-7.5) -- (0,4.5) node[left] {$y$};

        \node[below left] at (0,0) {$0$};
        \node[above] at (1,0) {$1$};
        \node[above left] at (0,1) {$1$};
        \node[below] at (5,0) {$5$};
        \node[below] at (-4,0) {$-4$};

        % Графік (хвиля)
        \draw[thick] plot [smooth, tension=0.44] coordinates {(-4, -6) (-3, -5) (-2, -3)(-1, 0) (0, 1)(1, 0) (2, -3) (3, -4) (5, -1)};
        \fill (-4,-6) circle (3pt);
        \fill (5, -1) circle (3pt);
        \node[below] at (3, -3.5) {\small $y=f(x)$};
    \end{tikzpicture}\end{nmtfit}
    \end{flushright}
\end{minipage}
\par\penalty-20\vspace{0.25cm}
\end{samepage}

\begin{samepage}
\noindent\zadnum Узгодьте функцію (1–3) із її властивістю (А – Д). \nmtyear{2025}

\nmtnobreak
\nopagebreak\nbvspace{0.3cm}

\nmtnobreak
\noindent
\matchingLayout{
\matchHead{Функція}
    \matchItem{1}{$y=\lg x$}
\matchItem{2}{$y=\dfrac{1}{x}$}
\matchItem{3}{$y=(x+1)^2$}
}{
\matchHead{Властивість функції}
    \matchItem{А}{зростає на всій області визначення}
\matchItem{Б}{парна}
\matchItem{В}{має точку максимуму}
\matchItem{Г}{не визначена лише в одній точці}
\matchItem{Д}{має точку мінімуму}
}{
\matchingGrid
}
\par\penalty-20\vspace{0.25cm}
\end{samepage}

\begin{samepage}
\noindent\zadnum На рисунку зображено графік функції $y=f(x)$, визначеної на проміжку $[-4; 5]$. Узгодьте проміжок (1–3) з властивістю (А – Д), яку має функція на цьому проміжку. \nmtyear{2025}

\nmtnobreak
\nopagebreak\nbvspace{0.2cm}

\nmtnobreak
\noindent
\begin{minipage}[t]{0.55\textwidth}
    \nopagebreak\nbvspace{0pt}
\matchHead{Проміжок}
    \matchItem{1}{$[-4; -2]$}
\matchItem{2}{$[-2; 2]$}
\matchItem{3}{$[2; 5]$}

    \nopagebreak\nbvspace{0.3cm}
\matchHead{Властивість функції}
    \matchItem{А}{функція спадає}
\matchItem{Б}{функція має точку максимуму}
\matchItem{В}{графік функції перетинає пряму $y=-4{,}5$}
\matchItem{Г}{для кожної точки $(x_0; y_0)$, що належить графіку функції, добуток $x_0 \cdot y_0 < 0$}
\matchItem{Д}{функція набуває лише додатних значень}

    \nopagebreak\nbvspace{0.3cm}

    \matchingGrid
\end{minipage}%
\hfill
\begin{minipage}[t]{0.40\textwidth}
    \nopagebreak\nbvspace{0pt}
    \begin{flushright}
    \begin{nmtfit}\begin{tikzpicture}[scale=0.5]
        \draw[step=1cm,gray!50,very thin] (-6.5,-7.5) grid (5.5,4.5);
        \draw[->, >=stealth, thick] (-6.5,0) -- (5.5,0) node[below] {$x$};
        \draw[->, >=stealth, thick] (0,-7.5) -- (0,4.5) node[left] {$y$};

        \node[below left] at (0,0) {$0$};
        \node[above] at (1,0) {$1$};
        \node[above left] at (0,1) {$1$};
        \node[below] at (5,0) {$5$};
        \node[below] at (-4,0) {$-4$};

        % Графік (хвиля)
        \draw[thick] plot [smooth, tension=0.44] coordinates {(-4, -6) (-3, -5) (-2, -3)(-1, 0) (0, 1)(1, 0) (2, -3) (3, -4) (5, -1)};
        \fill (-4,-6) circle (3pt);
        \fill (5, -1) circle (3pt);
        \node[below] at (3, -3.5) {\small $y=f(x)$};
    \end{tikzpicture}\end{nmtfit}
    \end{flushright}
\end{minipage}
\par\penalty-20\vspace{0.25cm}
\end{samepage}

\begin{samepage}
% НМТ 2026, сесія 23.05, завдання №18
\zadtask{Установіть відповідність між характеристикою функції \mbox{(1--3)} та проміжком \mbox{(А--Д)}, якому вона належить. \nmtyear{2026}

\nmtnobreak
\nopagebreak\nbvspace{0.3cm}

\nmtnobreak
\noindent
\matchingLayout{
\matchHead{Характеристика}
\matchItem{1}{область визначення функції $y = \sqrt{x - 2}$}
\matchItem{2}{точка екстремуму функції $y = x^2 + 2x + 3$}
\matchItem{3}{абсциса точки перетину прямої $y = 6$ з графіком функції $y = x^3$}
}{
\matchHead{Проміжок}
\matchItem{А}{$(-\infty;\, -2)$}
\matchItem{Б}{$[-2;\, 0)$}
\matchItem{В}{$[0;\, 1)$}
\matchItem{Г}{$(1;\, 2)$}
\matchItem{Д}{$[2;\, +\infty)$}
}{
\matchingGrid
}
}
% Відповідь: 1~--~Д, 2~--~Б, 3~--~Г
\ifshowsolutions\solution{$1)$ ОДЗ $y=\sqrt{x-2}$: $x\ge 2$, тобто $[2;+\infty)$~(Д). $2)$ Вершина параболи $y=x^2+2x+3$: $x=-\dfrac{2}{2}=-1\in[-2;0)$~(Б). $3)$ $x^3=6 \Rightarrow x=\sqrt[3]{6}\approx 1{,}8\in(1;2)$~(Г). Відповідь: \textbf{1--Д, 2--Б, 3--Г}.}\fi
\par\penalty-20
\end{samepage}

\begin{samepage}
% НМТ 2026, сесія 30.05, завдання №17
\zadtask{Установіть відповідність між характеристикою функції \mbox{(1--3)} та проміжком \mbox{(А--Д)}, яким є ця характеристика. \nmtyear{2026}

\nmtnobreak
\nopagebreak\nbvspace{0.2cm}
\noindent
\matchingLayout{
\matchHead{Характеристика функції}
\matchItem{1}{область значень функції\\ $y = \sqrt{x+2}$}
\matchItem{2}{область визначення функції \\$y = \lg(x-2)$}
\matchItem{3}{проміжок зростання функції\\ $y = -x^2 + 4x - 7$}
}{
\matchHead{Проміжок}
\matchItem{А}{$(-\infty;\, +\infty)$}
\matchItem{Б}{$(2;\, +\infty)$}
\matchItem{В}{$[0;\, +\infty)$}
\matchItem{Г}{$(-\infty;\, 2]$}
\matchItem{Д}{$(-\infty;\, 2)$}
}{
\matchingGrid
}
}
% Відповідь: 1~--~В, 2~--~Б, 3~--~Г
\ifshowsolutions\solution{\textbf{1}: область значень $y=\sqrt{x+2}$ — це $[0;+\infty)$ (В). \textbf{2}: ОДЗ $\lg(x-2)$: $x-2>0$, тобто $(2;+\infty)$ (Б). \textbf{3}: $y=-x^2+4x-7$ — парабола вітками вниз, вершина $x=2$, зростає на $(-\infty;2]$ (Г). Відповідь: \textbf{1~--~В, 2~--~Б, 3~--~Г}.}\fi
\par\penalty-20
\end{samepage}

\begin{samepage}
% НМТ 2026, сесія 1.06, завдання №17 --- основна тема: Графіки елементарних функцій, їх побудова та перетворення
\zadtask{Увідповідніть початок речення \mbox{(1--3)} із його закінченням \mbox{(А--Д)} так, щоб утворилося правильне твердження. \nmtyear{2026}

\nmtnobreak
\nopagebreak\nbvspace{0.2cm}
\noindent
\matchingLayout{
\matchHead{Початок речення}
\matchItem{1}{Графік функції $y = \dfrac{2}{x} + 3$}
\matchItem{2}{Графік функції $y = 4^x$}
\matchItem{3}{Графік функції $y = \mathrm{tg}\,x$}
}{
\matchHead{Закінчення речення}
\matchItem{А}{симетричний відносно початку координат.}
\matchItem{Б}{симетричний відносно осі $y$.}
\matchItem{В}{перетинає лише вісь $x$.}
\matchItem{Г}{перетинає лише вісь $y$.}
\matchItem{Д}{перетинає графік функції $y = x$ лише в одній точці.}
}{
\matchingGrid
}
}
% Відповідь: 1~--~В, 2~--~Г, 3~--~А
\ifshowsolutions\solution{\textbf{1)} $y=\dfrac{2}{x}+3$ ($x\ne0$): осі $y$ не перетинає, а вісь $x$ перетинає при $x=-\dfrac{2}{3}$ --- лише вісь $x$ (В); \textbf{2)} $y=4^x>0$: перетинає лише вісь $y$ у точці $(0;1)$ (Г); \textbf{3)} $y=\mathrm{tg}\,x$ --- непарна, симетрична відносно початку координат (А). Відповідь: \textbf{1\,--\,В, 2\,--\,Г, 3\,--\,А}.}\fi
\par\penalty-20
\end{samepage}

\begin{samepage}
% НМТ 2026, сесія 4.06, завдання №17
\zadtask{Установіть відповідність між функцією, заданою на проміжку \mbox{(1--3)}, та її найбільшим значенням \mbox{(А--Д)} на цьому проміжку. \nmtyear{2026}

\nmtnobreak
\nopagebreak\nbvspace{0.3cm}
\noindent
\matchingLayout{
\matchHead{Функція}
\matchItem{1}{$y = \sin 4x + 2,\ x \in (-\infty; +\infty)$}
\matchItem{2}{$y = 4 - x^2,\ x \in [0; 2]$}
\matchItem{3}{$y = 4 + \sqrt{x},\ x \in [0; 1]$}
}{
\matchHead{Найбільше значення}
\matchItem{А}{$0$}
\matchItem{Б}{$3$}
\matchItem{В}{$4$}
\matchItem{Г}{$5$}
\matchItem{Д}{$6$}
}{
\matchingGrid
}
}
% Відповідь: 1~--~Б, 2~--~В, 3~--~Г
\ifshowsolutions\solution{1: $\sin4x\leqslant1$, тож найбільше $y=1+2=3$ (Б). 2: $y=4-x^2$ спадає на $[0;2]$, найбільше при $x=0$: $y=4$ (В). 3: $y=4+\sqrt{x}$ зростає на $[0;1]$, найбільше при $x=1$: $y=5$ (Г). Відповідь: \textbf{1~--~Б, 2~--~В, 3~--~Г}.}\fi
\par\penalty-20
\end{samepage}

\begin{samepage}
% НМТ 2026, сесія 5.06, завдання №17
\zadnum Доберіть до функції \mbox{(1--3)} її властивість \mbox{(А--Д)}. \nmtyear{2026}

\nmtnobreak
\nopagebreak\nbvspace{0.3cm}
\noindent
\matchingLayout{
\matchHead{Функція}
\matchItem{1}{$y = x^4 - x^2$}
\matchItem{2}{$y = \log_{0{,}2} x$}
\matchItem{3}{$y = -\dfrac{1}{x}$}
}{
\matchHead{Властивість функції}
\matchItem{А}{область визначення збігається з областю значень}
\matchItem{Б}{є спадною}
\matchItem{В}{має лише три нулі}
\matchItem{Г}{має лише два нулі}
\matchItem{Д}{графік розташований у I та III чвертях}
}{
\matchingGrid
}
% Відповідь: 1~--~В, 2~--~Б, 3~--~А
\ifshowsolutions\solution{$1$: $y=x^4-x^2=x^2(x^2-1)$ має нулі $x=0;\,\pm1$, тобто три нулі (В); $2$: $y=\log_{0{,}2}x$ при основі $0<0{,}2<1$ спадна (Б); $3$: $y=-\dfrac1x$ має область визначення $x\ne0$, що збігається з областю значень (А). Відповідь: \textbf{1~--~В, 2~--~Б, 3~--~А}.}\fi
\par\penalty-20
\end{samepage}

\begin{samepage}
% НМТ 2026, сесія 8.06, завдання №17 --- також у темі: Показникова функція і показникові вирази
\zadtask{Доберіть до функції \mbox{(1--3)} її властивість \mbox{(А--Д)}. \nmtyear{2026}\par
\nopagebreak\nbvspace{0.3cm}
\noindent
\matchingLayout{
\matchHead{Функція}
\matchItem{1}{$y = \lg x$}
\matchItem{2}{$y = 0{,}2^x - 1$}
\matchItem{3}{$y = x^2 - 4x + 5$}
}{
\matchHead{Властивість}
\matchItem{А}{є зростаючою}
\matchItem{Б}{є спадною}
\matchItem{В}{не має нулів}
\matchItem{Г}{має лише два нулі}
\matchItem{Д}{є парною}
}{
\matchingGrid
}
}
\par\penalty-20
\end{samepage}

\begin{samepage}
% НМТ 2026, сесія 9.06, завдання №17
\zadtask{Доберіть до функції \mbox{(1--3)} її властивість \mbox{(А--Д)}. \nmtyear{2026}\par
\nopagebreak\nbvspace{0.3cm}
\noindent
\matchingLayout{
\matchHead{Функція}
\matchItem{1}{$y = \sin x$}
\matchItem{2}{$y = x^2 + 2$}
\matchItem{3}{$y = 3x + 2$}
}{
\matchHead{Властивість функції}
\matchItem{А}{парна}
\matchItem{Б}{графік функції симетричний відносно осі $x$}
\matchItem{В}{спадає на області визначення}
\matchItem{Г}{зростає на області визначення}
\matchItem{Д}{графік функції проходить через точку $(0; 0)$}
}{
\matchingGrid
}
}
\par\penalty-20
\end{samepage}

\begin{samepage}
% НМТ 2026, сесія 10.06, завдання №17 --- основна тема: Похідна
\zadtask{На рисунку зображено графік функції $y = f(x)$, визначеної на проміжку $[-2;\, 4]$. Для будь-якого значення $x$ із проміжку $(-2;\, 4)$ існує похідна $f'(x)$. До кожного початку \mbox{(1--3)} речення доберіть його закінчення \mbox{(А--Д)} так, щоб утворилося правильне твердження. \nmtyear{2026}\par
\nopagebreak\nbvspace{0.3cm}
\begin{center}
\begin{nmtfit}\begin{tikzpicture}[scale=0.55, thick]
    \draw[step=1cm, gray!40, thin] (-3,-4) grid (5,6);
    \draw[->] (-3.5,0) -- (5.5,0) node[right] {$x$};
    \draw[->] (0,-4.5) -- (0,6.5) node[above] {$y$};
    \node[below left, font=\small] at (0,0) {$0$};
    \node[below right, font=\small] at (1,0) {$1$};
    \node[above left, font=\small] at (0,1) {$1$};
    \node[left, font=\small] at (0,5) {$5$};
    \node[below, font=\small] at (-2,0) {$-2$};
    \node[below, font=\small] at (4,0) {$4$};
    \node[right, font=\small] at (0.15,-3) {$-3$};
    \draw[magenta!80!black, thick, smooth] plot coordinates {(-2,-3) (-1.3,0.3) (-0.6,1.7) (0,2) (0.7,1.5) (1.4,1.1) (2,1) (2.7,1.6) (3.3,3) (4,5)};
    \fill (-2,-3) circle (3pt);
    \fill (4,5) circle (3pt);
    \node[right, font=\small, fill=white, inner sep=1.5pt] at (1.7,3.6) {$y=f(x)$};
\end{tikzpicture}\end{nmtfit}
\end{center}
\nopagebreak\nbvspace{0.2cm}
\noindent
\matchingLayout{
\matchHead{Початок речення}
\matchItem{1}{Найбільше значення функції $y=f(x)$ на проміжку $[-2;\, 4]$ дорівнює}
\matchItem{2}{Функція набуває від'ємного значення, якщо $x =$}
\matchItem{3}{$f'(x) = 0$, якщо $x =$}
}{
\matchHead{Закінчення речення}
\matchItem{А}{$-2$.}
\matchItem{Б}{$-1$.}
\matchItem{В}{$2$.}
\matchItem{Г}{$4$.}
\matchItem{Д}{$5$.}
}{
\matchingGrid
}
}
\par\penalty-20
\end{samepage}

\begin{samepage}
% НМТ 2026, сесія 12.06, завдання №17 --- також у темі: Показникова функція і показникові вирази
\zadtask{Доберіть до функції \mbox{(1--3)} її властивість \mbox{(А--Д)}. \nmtyear{2026}\par
\nopagebreak\nbvspace{0.3cm}
\noindent
\matchingLayout{
\matchHead{Функція}
\matchItem{1}{$y = \sqrt{x - 3}$}
\matchItem{2}{$y = x^2 + 3$}
\matchItem{3}{$y = 3^x$}
}{
\matchHead{Властивість функції}
\matchItem{А}{областю значень є проміжок $[3; +\infty)$}
\matchItem{Б}{областю визначення є проміжок $[3; +\infty)$}
\matchItem{В}{областю визначення є проміжок $[-3; +\infty)$}
\matchItem{Г}{зростає на множині всіх дійсних чисел}
\matchItem{Д}{є непарною}
}{
\matchingGrid
}
}
\par\penalty-20
\end{samepage}

\begin{samepage}
% НМТ 2026, сесія 16.06, завдання №17 --- основна тема: Графіки елементарних функцій, їх побудова та перетворення
\zadtask{Доберіть до функції \mbox{(1--3)} властивість її графіка \mbox{(А--Д)}. \nmtyear{2026}\par
\nopagebreak\nbvspace{0.3cm}
\noindent
\matchingLayout{
\matchHead{Функція}
\matchItem{1}{$y = \sqrt{x-2} + 1$}
\matchItem{2}{$y = -5^x$}
\matchItem{3}{$y = x^3$}
}{
\matchHead{Властивість графіка функції}
\matchItem{А}{симетричний відносно початку координат}
\matchItem{Б}{симетричний відносно осі $y$}
\matchItem{В}{розміщений лише в ІІІ та IV координатних чвертях}
\matchItem{Г}{розміщений лише в І координатній чверті}
\matchItem{Д}{проходить через точку $(0; 1)$}
}{
\matchingGrid
}
}
% Відповідь: 1--Г; 2--В; 3--А
\par\penalty-20
\end{samepage}

\typeTitle{Завдання з короткою відповіддю}

\begin{samepage}
\noindent\zadnum Функція $f(x)$ набуває значення $3$ для всіх додатних значень $x$ і значення $1$ — для всіх від'ємних значень $x$. Функція $y=g(x)$ є прямою пропорційністю, а її графік проходить через точку $(1; 0{,}5)$. Обчисліть значення виразу $\dfrac{5g(-1) + 3f(2)}{f(-1)}$. \nmtyear{2025}

\nmtnobreak
\nopagebreak\nbvspace{0.3cm}
\noindent\textbf{Відповідь: \underline{\hspace{3cm}}}
\nmtAnswerBox
\par\penalty-20
\end{samepage}

\begin{samepage}
\noindent\zadnum Функція $y=f(x)$ є прямою пропорційністю на проміжку $(-\infty; +\infty)$. Знайдіть значення функції $g(x) = 3f(x) + \dfrac{|x|}{8}$ у точці з абсцисою $x_0 = -2$, якщо графік функції $y=f(x)$ проходить через точку $(2; 7)$. \nmtyear{2025}

\nmtnobreak
\nopagebreak\nbvspace{0.3cm}
\noindent\textbf{Відповідь: \underline{\hspace{3cm}}}
\nmtAnswerBox
\par\penalty-20
\end{samepage}

\begin{samepage}
\noindent\zadnum Функція $f(x)$ є парною, а функція $g(x)$ — непарна. Обчисліть значення виразу $\dfrac{4g(-x_0) + 2f(x_0)}{5f(-x_0)}$, якщо $f(x_0)=1$, а $g(x_0)=-3$. \nmtyear{2025}

\nmtnobreak
\nopagebreak\nbvspace{0.3cm}
\noindent\textbf{Відповідь: \underline{\hspace{3cm}}}
\nmtAnswerBox
\par\penalty-20
\end{samepage}

\ifshowsolutions
\sectionTitle{ВІДПОВІДІ ДО ЗАВДАНЬ НМТ 2026}
\begin{multicols}{3}\noindent
\noindent\textbf{20.}~Д \ {\scriptsize\color{gray}(2.06, №8)}\par
\noindent\textbf{21.}~В \ {\scriptsize\color{gray}(3.06, №8)}\par
\noindent\textbf{22.}~В \ {\scriptsize\color{gray}(4.06, №8)}\par
\noindent\textbf{23.}~А \ {\scriptsize\color{gray}(5.06, №8)}\par
\noindent\textbf{26.}~Г \ {\scriptsize\color{gray}(15.06, №8)}\par
\noindent\textbf{27.}~А \ {\scriptsize\color{gray}(16.06, №8)}\par
\noindent\textbf{28.}~В \ {\scriptsize\color{gray}(17.06, №8)}\par
\noindent\textbf{29.}~А \ {\scriptsize\color{gray}(19.06, №8)}\par
\noindent\textbf{51.}~1~--~Д, 2~--~Б, 3~--~Г \ {\scriptsize\color{gray}(23.05, №18)}\par
\noindent\textbf{52.}~1~--~В, 2~--~Б, 3~--~Г \ {\scriptsize\color{gray}(30.05, №17)}\par
\noindent\textbf{53.}~1~--~В, 2~--~Г, 3~--~А \ {\scriptsize\color{gray}(1.06, №17)}\par
\noindent\textbf{54.}~1~--~Б, 2~--~В, 3~--~Г \ {\scriptsize\color{gray}(4.06, №17)}\par
\noindent\textbf{55.}~1~--~В, 2~--~Б, 3~--~А \ {\scriptsize\color{gray}(5.06, №17)}\par
\noindent\textbf{60.}~1--Г; 2--В; 3--А \ {\scriptsize\color{gray}(16.06, №17)}\par
\end{multicols}
\fi

\endgroup
\end{document}
